% =============================================================================
% Empirical Facts section for model.tex
% Insert with: % =============================================================================
% Empirical Facts section for model.tex
% Insert with: % =============================================================================
% Empirical Facts section for model.tex
% Insert with: % =============================================================================
% Empirical Facts section for model.tex
% Insert with: \input{empirical_facts} between \section{Introduction} and
% \section{The Model}, i.e. around line 268 of model.tex.
% =============================================================================

\section{Empirical Facts} \label{sec:empirical_facts}

This section presents four empirical regularities that motivate the model and
anchor its primitives in the Indian case. Subsection~\ref{subsec:d_narrative}
documents the staggered intensification of the symbolic-policy instrument~$d$
and the co-movement of majority and minority narrative mobilization,
$E_M$~and~$E_m$. Subsection~\ref{subsec:cycles} shows the co-existence of
persistent Hawk-state absorption and recurrent regime cycles across Indian
states, the two dynamic paths identified by Propositions~\ref{prop:trap}
and~\ref{prop:cycles}. Subsection~\ref{subsec:voter_space} places states in the
two-dimensional space $(\Delta u,\text{BJP dominance})$ with ban intensity as a
third margin, tracing the empirical footprint of the Security Premium of
Lemma~\ref{lem:reaction_invertedU}. Subsection~\ref{subsec:other_democracies}
closes with three scope cases --- the United States, Colombia, and Iran ---
that illustrate how the same mechanism operates under different institutional
architectures. The full set of data sources, variable constructions, and
supplementary results is deferred to Appendix~\ref{app:more_empirical_facts}.

\subsection{Symbolic Policy and Narrative Mobilization}
\label{subsec:d_narrative}

India's post-independence polity inherited a constitutional encouragement of
cow protection (Article~48) but no national prohibition. Through the first
four post-independence decades, beef regulation remained a state-level matter,
with Congress-governed states tolerating buffalo and low-yield cattle
slaughter and a handful of Jan Sangh- or BJP-governed states pushing stricter
enforcement. The legal landscape changed discontinuously after the BJP's
national victory in 2014. Between 2015 and 2017, a wave of state-level
legislation extended the ban from cows to the entire bovine family, reversed
the burden of proof, and raised maximum penalties to up to ten years'
imprisonment.\footnote{Maharashtra Animal Preservation (Amendment) Act, 2015;
Haryana Gauvansh Sanrakshan Act, 2015; analogous amendments in Gujarat,
Jharkhand, Uttar Pradesh, Madhya Pradesh, and Uttarakhand. A 2017 Ministry of
Environment notification sought to restrict cattle sales nationally before
being stayed by the Supreme Court.} \autoref{fig:ban_adoption} records the
cumulative adoption of six ban types across states, and Map~A
(\autoref{fig:maps}) summarizes cumulative ban intensity as of 2020.

\textbf{Majority-side narrative supply.} To trace the co-movement of the
symbolic instrument $d$ with majority-side mobilization $E_M$, I assemble a
corpus of $2{,}324$ BJP leadership texts from three sources: the party's
official speech archive
(\href{https://www.bjp.org/speeches}{bjp.org/speeches}, 1996--2025), the Prime
Minister's office
(\href{https://www.pmindia.gov.in}{pmindia.gov.in}, 2023--2026), and BJP
Members' floor speeches in the Lok Sabha, harvested from the official DSpace
backend of \href{https://elibrary.sansad.in}{elibrary.sansad.in}
(1996--2024). The Lok Sabha supplement fills the pre-2014 gap in the party's
online archive and captures parliamentary, as distinct from campaign,
register. Speakers include the Prime Minister, the four BJP National
Presidents of the period (Amit Shah, Rajnath Singh, J.~P.~Nadda, and Narendra
Modi himself), and senior front-benchers (Arun Jaitley, Sushma Swaraj, Nitin
Gadkari). Each speech is classified into one of four categories
by a keyword-lexicon procedure: (i)~\emph{identity-coded} content
(cow/beef/gau, Hindu--Muslim pairs, love jihad, CAA/NRC, ghar wapsi, Ayodhya,
UCC, hijab, loudspeaker, Gyanvapi, waqf, triple talaq, and cognate terms in
Hindi); (ii)~\emph{economic-programmatic} content (MSME, PM-KISAN, Ujjwala,
Aadhaar-linked DBT, GST, ``Make in India,'' ``Viksit Bharat,'' and cognates);
(iii)~\emph{security-and-foreign} content (Pakistan, Galwan, Article~370,
Operation Sindoor); (iv)~\emph{residual}.

\textbf{Minority-side narrative.} A symmetric majority-side corpus proved
infeasible: of the three major umbrella organizations --- the All India
Majlis-e-Ittehadul Muslimeen (AIMIM), the Jamiat Ulama-i-Hind, and the All
India Muslim Personal Law Board --- none operates a crawlable online speech
archive comparable to \texttt{bjp.org}.\footnote{As of April~2026,
\texttt{aimim.in} fails to resolve; \texttt{jamiatulama.org} resolves to a
parked GoDaddy lander; \texttt{aimplboard.in} retains a static early-2000s
site whose press-release endpoint returns~500. The archives of Owaisi's public
speeches persist on social-media platforms but are not systematically
catalogued with dates and transcripts. This asymmetry in narrative
infrastructure is itself informative: it is the observable expression of the
organizational-capacity parameter $\kappa$ in the Minority's collective-action
problem (\S\ref{subsec:joint_contest}), and tightens rather than weakens the
model's claim that the Minority's mobilization is supply-constrained.} I
therefore construct a minority-side salience series from a proxy corpus of
$3{,}783$ dated Indian news items, $2010$--$2025$, curated from major
English-language outlets --- The Hindu, The Times of India, NDTV, India Today,
Indian Express, Al Jazeera, Scroll, Maktoob, SabrangIndia, Muslim Mirror, and
the BBC --- using $65$ keyword queries covering AIMIM/Owaisi, Jamiat, AIMPLB,
cow-related lynchings (Akhlaq, Pehlu Khan, Junaid, Tabrez), Shaheen Bagh,
CAA/NRC, hijab, waqf, UCC, Gyanvapi, Babri, Bilkis, and Kashmir Muslim
leadership. Each item is classified into identity-coded, counter-mobilization
(protest, bandh, satyagraha, rally, ``not in my name,'' boycott),
security-and-foreign, economic, or residual.

\textbf{Results.} \autoref{fig:narrative_timeseries} plots the monthly share
of identity-coded content on the majority side (top panel) and the combined
identity-coded and counter-mobilization share on the minority side (bottom
panel), smoothed over a three-month rolling window. Four patterns are
visible. First, the majority-side identity share steps up in 2015--17 around
the state-level ban wave, contracts during the Covid period, and re-expands
in 2022--24 around the Uniform Civil Code and Waqf debates. Second, the
minority-side share is quiescent in 2010--14, rises discretely with the 2015
Dadri lynching and the 2017 Pehlu Khan killing, peaks during the 2019--20
CAA/NRC and Shaheen Bagh episode, and remains elevated through 2025. Third,
the two series move together at high frequency around legislative shocks
(2015 state bans, 2019 CAA), consistent with the joint contest structure of
Lemma~\ref{lem:contest_equilibrium} in which $d>0$ activates both
$E_M^{*}$~and~$E_m^{*}$. Fourth, the episode-level evidence is consistent
with the inverted-U bound on the Security Premium
(Lemma~\ref{lem:reaction_invertedU}): the majority-side identity share peaks
in 2016--17 and again in 2024, with intermediate compressions that coincide
with electorally costly overreach events (the 2018 Supreme Court stay of the
livestock-market notification; the 2020--21 farm laws, which were ultimately
repealed).

% ============================================================================
\begin{figure}[t]
    \centering
    \includegraphics[width=0.92\textwidth]{figures/empirics/narrative_timeseries.pdf}
    \caption{\normalsize\bf Majority and Minority Identity-Coded Narrative
    Share, India 2010--2025.}
    \floatfoot*{\footnotesize\textit{Notes.} Top panel: share of BJP
    leadership texts classified as identity-coded in the month,
    $N=2{,}324$ from \texttt{bjp.org/speeches}, \texttt{pmindia.gov.in}, and
    BJP Lok Sabha floor speeches harvested from
    \texttt{elibrary.sansad.in}, 1996m9--2026m4. Bottom panel: share of
    curated Muslim-voice items classified as identity-coded or
    counter-mobilization in the month, $N=3{,}792$ items
    ($3{,}783$~news-outlet proxies plus~$9$ AIMIM Lok~Sabha floor
    speeches), 2004m12--2025m12. Both series are three-month rolling
    means. Vertical lines mark the Dadri lynching (Sep~2015), the Haryana
    Gauvansh Sanrakshan Act (Mar~2015), the Pehlu Khan lynching (Apr~2017),
    the Citizenship (Amendment) Act passage and onset of Shaheen Bagh
    (Dec~2019), the farm laws (Sep~2020) and their repeal (Nov~2021), the
    Mahsa Amini episode in Iran (Sep~2022), and the 2024 Kanwar Yatra. The
    classification lexicon and methodology are reported in
    Appendix~\ref{app:narrative_classification}. Sources: bjp.org/speeches;
    pmindia.gov.in/en/news\_updates; curated news-outlet RSS archives.}
    \label{fig:narrative_timeseries}
\end{figure}
% ============================================================================

% ============================================================================
\begin{figure}[t]
    \centering
    \includegraphics[width=0.82\textwidth]{figures/empirics/map_ban_intensity.pdf}
    \caption{\normalsize\bf Composite Cattle-Slaughter Ban Intensity by
    State, Cumulative Through~2020.}
    \floatfoot*{\footnotesize\textit{Notes.} Composite ban intensity is
    computed as the sum over six ban types (cow, bullock, bull, buffalo
    slaughter; beef possession; beef sales) of
    $\max(\text{current year}-\text{adoption year},0)$, averaged over
    1990--2020. Higher values indicate earlier and broader adoption. Shapefile:
    GADM~v4.1 administrative level~1. Underlying policy
    data: Cattle Slaughter Policy dataset, constructed from state-level
    legislative records; see Appendix~\ref{app:ban_intensity}.}
    \label{fig:ban_adoption}
\end{figure}
% ============================================================================

\subsection{Cycles and Persistence in India's Political Competition}
\label{subsec:cycles}

The dynamic extension of the model in Section~\ref{sec:dynamic} delivers two
qualitatively distinct political paths: with sufficiently patient voters,
the Hawk becomes an absorbing state of the polity (Proposition~\ref{prop:trap});
with myopic retrospective voters, the same mechanism generates recurrent
Hawk--Dove regime alternation (Proposition~\ref{prop:cycles}). India's
federal structure provides 30 quasi-independent state polities observed over
1984--2024, in which both paths can be distinguished simultaneously.

I construct a state-year panel of Vidhan Sabha (state legislative assembly)
outcomes by combining the Bhavnani India Election Dataset v2.0 (1977--2015)
with supplementary state-year records I harvested from the Election
Commission of India and Wikipedia for post-2015 elections (the $51$ Vidhan
Sabha elections and the 2019 and 2024 Lok Sabha elections held in the
extension window). The panel is forward-filled between elections, so each
state-year is assigned the incumbent party of the preceding Assembly
result. Using the BJP's seat share, the top-1 party's vote share margin, and
the incumbent party identity, I classify each state-year as
\emph{Hawk-dominant} if the BJP rules with either a seat-share majority or a
vote-share margin of at least~$10$~pp;
\emph{Dove-dominant} if a non-BJP party rules with a margin of at
least~$10$~pp; and \emph{competitive} otherwise.

\autoref{fig:cycles_four_panel} presents four panels derived from the panel.
Panel~(A) plots the national BJP vote and seat shares in Lok Sabha elections
from~1984 to~2024. Two step increases stand out: the 1999--2004 NDA years
and the post-2014 Modi wave, with a peak of~55.8\%~of seats in~2019 on a
$37.4\%$~vote share. Panel~(B) plots the fraction of state Assemblies
governed by the BJP across time: the share rises from near-zero before~1990
to~$0.55$ after~2017. Panel~(C) plots the distribution of uninterrupted
incumbency tenure at the state level, by party group. The BJP distribution
has a pronounced right tail: the party holds five state spells of at
least~11~consecutive years, including Gujarat 1995--2024 (30~years),
Madhya~Pradesh 2003--2018 (15~years), Chhattisgarh 2003--2018 (15~years), and
Uttar~Pradesh 1991--2001 (11~years). The INC and regional-party
distributions are centred on shorter tenures. Panel~(D) plots the state-year
regime heatmap. Three patterns are visible. First, the post-2014 Hindi Belt
--- U.P., M.P., Rajasthan, Chhattisgarh, Gujarat, Haryana, Jharkhand ---
is populated by long Hawk-dominant spells with few transitions, the
empirical footprint of the conflict trap. Second, the southern and eastern
states --- Tamil~Nadu, Kerala, West~Bengal, Andhra~Pradesh, Telangana ---
remain Dove-dominant throughout, consistent with $\delta>\bar{\delta}$ not
applying uniformly. Third, an intermediate set of states --- Himachal~Pradesh
(seven regime transitions), Rajasthan ($6$), Haryana ($5$), Karnataka~($4$),
and Delhi --- alternate between categories, the empirical footprint of the
cycles case.

% ============================================================================
\begin{figure}[t]
    \centering
    \includegraphics[width=0.98\textwidth]{figures/empirics/cycles_four_panel.pdf}
    \caption{\normalsize\bf Cycles and Persistence in Indian Electoral
    Competition, 1984--2024.}
    \floatfoot*{\footnotesize\textit{Notes.} (A)~National BJP vote share and
    Lok Sabha seat share, 1984--2024. (B)~Fraction of 30~Indian states (and
    NCT of Delhi) governed by the BJP or a BJP-led coalition in each year.
    (C)~Histogram of uninterrupted Assembly-level incumbency tenure by party
    group \{BJP, INC, regional/other\}; spells that are still ongoing at the
    end of the panel are right-truncated. (D)~State-year regime classification
    (rows: $30$ states, alphabetical; columns: 1984--2024). A state-year is
    \emph{Hawk-dominant} when the BJP rules with seat-share
    $\geq 50\%$~or~vote-share margin~$\geq 10$~pp;
    \emph{Dove-dominant} when a non-BJP party rules with margin
    $\geq 10$~pp; \emph{competitive} otherwise.
    Data: Bhavnani India Election Dataset v2.0 (1977--2015) and
    author-compiled supplementary panel from Wikipedia and ECI for the
    $51$~Vidhan Sabha elections and 2019, 2024 Lok Sabha elections
    held in 2015--2025. See Appendix~\ref{app:data_sources} for the
    construction and sanity checks.}
    \label{fig:cycles_four_panel}
\end{figure}
% ============================================================================

The co-existence of long Hawk spells and intermediate-state cycles within the
same federation is not, by itself, a test of
Propositions~\ref{prop:trap} or~\ref{prop:cycles}: the classification is
reduced-form, and cross-state heterogeneity in~$\delta$, in the moderate-
defection technology $\chi_{\mathrm{mod}}(d)$, and in the primitive
state-capacity gap $S_A/S_B$ will all push states toward different paths.
What the panel does rule out is a uniform story --- either global
entrenchment or global alternation --- and is consistent with the model's
joint prediction that the same mechanism can produce both outcomes along
different parameter paths. Gujarat's 30-year Hawk spell, in particular, is the
archetype of the conflict trap: a continuous stretch in which the BJP
retained power through the 2002 communal violence and its aftermath, three
Lok Sabha electoral cycles, and at least one episode of sharp economic
contraction (the 2002--03 slowdown). That it did so while retaining
Dove-dominant neighbours (Rajasthan pre-2013, the 2018 swing) is exactly the
cross-state heterogeneity the model predicts.

\subsection{Voter Welfare and Political Dominance}
\label{subsec:voter_space}

Proposition~\ref{prop:provocation} states that the Hawk's electoral survival
requires turning the joint contest on when~$\Delta u>0$, and that the
symbolic-policy path is chosen to substitute an activated Security Premium
for the Hawk's economic deficit. The empirical footprint of this
substitution is a state-year distribution in the $(\Delta u,\text{BJP
dominance})$~space in which states with high cumulative ban intensity cluster
in a region distinct from states that have not used symbolic policy. This
subsection marks that distribution.

\textbf{Measurement.} I proxy $\Delta u$ by state GDP per capita relative to
the national average (where $1.0$ denotes the national average in the given
year), taken from the Reserve Bank of India's relative per-capita income
series at $1990, 2000, 2010, 2020$, and~$2023$ anchor years. I proxy BJP
dominance by the BJP's Vidhan Sabha seat share in the state-year, computed
from the election panel of \S\ref{subsec:cycles}. Composite ban intensity is
the state-level construction of \S\ref{subsec:d_narrative} carried forward
through~2020.

\textbf{Maps.} Maps~B and~C in \autoref{fig:maps} visualize the cross-section.
Map~A (ban intensity, already shown as \autoref{fig:ban_adoption}) clusters
on the northern and western states: Gujarat, Haryana, Punjab, Himachal
Pradesh, Uttarakhand, and large parts of Uttar Pradesh and Bihar. Map~B
plots Hindu--Muslim violence intensity per capita, pooling incidents from
the Varshney--Wilkinson dataset ($1950$--$1995$) and the GDELT Hindu--Muslim
variant ($1980$--$2024$), normalized by state population.\footnote{State
population is interpolated annually from Census~2011 shares rescaled by the
World~Bank national-population series; see
Appendix~\ref{app:population}.} The violence map co-locates with the ban map
along the Gangetic-plain belt and Gujarat. Map~C plots the change in BJP
state-Assembly vote share between the 2009-anchor and the 2019-anchor Vidhan
Sabha elections. The post-2014 breakthrough states ---Tripura (+42~pp),
Manipur (+35~pp), West~Bengal (+34~pp), Haryana (+27~pp), Uttar~Pradesh
(+23~pp), Assam (+22~pp), Odisha (+17~pp) --- are precisely the states that
had not previously delivered a BJP government and in which the post-2014 wave
represents an entry rather than a consolidation.

\textbf{Two-dimensional scatter.}
\autoref{fig:voter_space_2d} places each state-year in the $(\Delta
u,\text{BJP seat share})$ space, with ban intensity mapped to point size and
year to point color, for the anchor years $1990, 2000, 2010, 2020, 2023$.
Two centroids are marked. For the high-ban-intensity states (top tercile),
the $2000$~centroid sits at $(1.21, 0.27)$ --- above average relative income
and appreciable BJP dominance --- and the $2020$~centroid sits at
$(0.79, 0.33)$: relative income has fallen below the national average
while BJP dominance has risen. For the low-ban-intensity states (bottom
tercile), the $2000$~centroid sits at $(0.77, 0.01)$ and the $2020$~centroid
at $(0.41, 0.02)$: relative income has fallen sharply, while BJP dominance
is negligible throughout. The two arrows together are the empirical
signature of Proposition~\ref{prop:provocation}: in states where symbolic
policy was used intensively, political dominance stayed high (rose, in
fact) even as the relative-income position deteriorated; in states where
symbolic policy was not used, the collapse in relative income was not
offset. The relative movement in relative income is larger than the
movement in BJP dominance, consistent with the provocation-ceiling bound of
equation~\eqref{eq:delta_u_bar_bound}: the electoral offset provided by
activated symbolic contest is structurally limited, and even intense use of
$d$ cannot fully compensate for large income deficits.

% ============================================================================
\begin{figure}[t]
    \centering
    \subfigure[Composite cattle-slaughter ban intensity (cumulative through~2020).]{
        \includegraphics[width=0.32\textwidth]{figures/empirics/map_ban_intensity.pdf}
        \label{fig:map_ban_intensity}
    }
    \subfigure[Hindu--Muslim violence intensity per capita (1950--2024).]{
        \includegraphics[width=0.32\textwidth]{figures/empirics/map_conflict_intensity.pdf}
        \label{fig:map_conflict_intensity}
    }
    \subfigure[$\Delta$~BJP Vidhan Sabha vote share, 2009~$\to$~2019 (pp).]{
        \includegraphics[width=0.32\textwidth]{figures/empirics/map_bjp_delta_2009_2019.pdf}
        \label{fig:map_bjp_delta}
    }
    \caption{\normalsize\bf Spatial Distribution of Symbolic-Policy Intensity,
    Hindu--Muslim Violence, and BJP Breakthrough.}
    \floatfoot*{\footnotesize\textit{Notes.} (A)~Composite ban intensity as
    defined in \autoref{fig:ban_adoption}. (B)~Pooled Hindu--Muslim violent
    incidents per capita, sources: Varshney--Wilkinson 1950--1995 and GDELT
    Hindu--Muslim variant 1980--2024; population panel at
    Appendix~\ref{app:population}. (C)~BJP state-Assembly vote-share change
    between the Vidhan Sabha election closest to 2009 and the Vidhan Sabha
    election closest to 2019 for each state, in percentage points;
    diverging colour scale centred at zero.}
    \label{fig:maps}
\end{figure}
% ============================================================================

% ============================================================================
\begin{figure}[t]
    \centering
    \includegraphics[width=0.88\textwidth]{figures/empirics/voter_space_2d.pdf}
    \caption{\normalsize\bf Voter Welfare and Political Dominance, Indian
    States~1990--2023.}
    \floatfoot*{\footnotesize\textit{Notes.} Each point is a state-year
    observation. Horizontal axis: state GDP per capita relative to the
    national average (1.0~=~national average in the year); vertical axis:
    BJP state-Assembly seat share; point size: composite cumulative ban
    intensity (see \autoref{fig:ban_adoption}); point colour: year. Anchor
    years $\{1990, 2000, 2010, 2020, 2023\}$. Two arrows track the centroid
    movement of the top-tercile and bottom-tercile ban-intensity states
    between 2000 and 2020: high-ban states moved from $(1.21, 0.27)$ to
    $(0.79, 0.33)$; low-ban states moved from $(0.77, 0.01)$ to
    $(0.41, 0.02)$. The five states with the largest post-2014
    seat-share movement are labelled.}
    \label{fig:voter_space_2d}
\end{figure}
% ============================================================================

\subsection{Other Electoral Democracies: US, Colombia, and Iran}
\label{subsec:other_democracies}

The model is presented in the context of India's cow-slaughter and citizenship
legislation, but its primitives are generic to two-party electoral environments
with a salient out-group. Three further cases illustrate the mechanism
under different institutional architectures.

\emph{The United States, post-2015.}
The Republican Party's post-2015 turn to immigration enforcement and
culturally coded domestic security fits the model closely. The Democratic
comparative advantage lies in universalist economic provision (the Affordable
Care Act, the American Rescue Plan, the Inflation Reduction Act), consistent
with the Dove's role in $(S_B,\bar g_B)$. The symbolic instruments have been
legible: Executive Order~13769, border-wall construction, the 2018
family-separation policy, and the federalisation of culture-war education
policy. Each is a $d$-type policy with no direct $g$ payoff. The 2018
midterms, held immediately after the family-separation episode, returned a
Democratic House majority, consistent with
Lemma~\ref{lem:reaction_invertedU}'s inverted-U prediction: beyond
$d^{\mathrm{peak}}$, the Hawk's moderate-defection cost $\chi_{\mathrm{mod}}$
accumulates faster than the Security Premium. The dynamic path since~2016
---alternation~2016/2020/2024 --- is closer to the cycles case of
Proposition~\ref{prop:cycles} than to the trap of
Proposition~\ref{prop:trap}, consistent with shorter voter horizons and with
the institutional absence of the candidacy-gatekeeping that supports the
Indian and Iranian trap paths.

\emph{Colombia, the 2016 peace plebiscite.}
Colombia's post-2002 Uribista period extended the mechanism to an environment
in which the Minority is an ideologically coded armed party (the FARC) rather
than an ethnic or religious domestic out-group. The \emph{Democratic
Security} doctrine of Álvaro Uribe (2002--2010) and its continuation in the
Centro Democrático after~2014 fit the Hawk type: it promised visible coercive
capacity against FARC, imposed fiscal costs (the ``democratic security
tax''), and activated a two-sided contest in which paramilitary mobilisation
rose alongside state enforcement. The 2016 plebiscite on the Santos--FARC
peace accord is the sharpest case of Lemma~\ref{lem:dove_corner} cutting in
reverse: when the Dove attempted to move the system to $d=0$, enough voters
preferred the Hawk's manufactured-threat environment to block the transition.
That the successor administration subsequently implemented substantial
portions of the same accord demonstrates that the narrow plebiscite margin
was the binding electoral constraint, not a technical defect in the agreement.

\emph{Iran, morality policing and the 2022 Amini episode.}
The Islamic Republic of Iran adapts the model to a hybrid regime in which
competitive elections occur within candidacy-gatekeeping institutions (the
Guardian Council). The hard-line bloc associated with the Islamic
Revolutionary Guard Corps occupies the Hawk type; its most visible symbolic
instrument is the regulation of women's public conduct --- compulsory hijab,
the \emph{gasht-e ershad} morality police, and periodic ``bad hijab''
campaigns. The Mahsa Amini episode of September~2022 and the ensuing
\emph{Woman, Life, Freedom} protests are a textbook instance of the
joint-contest structure of Lemma~\ref{lem:contest_equilibrium}: a $d$~shock
triggered a large $E_m^*$, which in turn activated the regime's $S$~apparatus
(the Basij). The re-intensification of morality-police enforcement after the
protest's decline is the trap of Proposition~\ref{prop:trap} re-closing ---
but with a specifically institutional interpretation: the ``patience'' that
sustains the Hawk in Iran is supplied by candidacy gatekeeping rather than
by voter preferences over permanence, as the model's reading of
candidacy-filtered regimes anticipates.

Together, the three cases show that the India-centric analysis of
\S\S\ref{subsec:d_narrative}--\ref{subsec:voter_space} is an instance of a
more general pattern rather than a country-specific artefact. In each case
the incumbent's symbolic-policy instrument substitutes an activated contest
environment for a structural economic disadvantage, the inverted-U bound
disciplines the use of the instrument, and the dynamic path is shaped by the
interaction of voter patience, institutional gatekeeping, and the strength
of the moderate-defection channel.

% End of \section{Empirical Facts}
 between \section{Introduction} and
% \section{The Model}, i.e. around line 268 of model.tex.
% =============================================================================

\section{Empirical Facts} \label{sec:empirical_facts}

This section presents four empirical regularities that motivate the model and
anchor its primitives in the Indian case. Subsection~\ref{subsec:d_narrative}
documents the staggered intensification of the symbolic-policy instrument~$d$
and the co-movement of majority and minority narrative mobilization,
$E_M$~and~$E_m$. Subsection~\ref{subsec:cycles} shows the co-existence of
persistent Hawk-state absorption and recurrent regime cycles across Indian
states, the two dynamic paths identified by Propositions~\ref{prop:trap}
and~\ref{prop:cycles}. Subsection~\ref{subsec:voter_space} places states in the
two-dimensional space $(\Delta u,\text{BJP dominance})$ with ban intensity as a
third margin, tracing the empirical footprint of the Security Premium of
Lemma~\ref{lem:reaction_invertedU}. Subsection~\ref{subsec:other_democracies}
closes with three scope cases --- the United States, Colombia, and Iran ---
that illustrate how the same mechanism operates under different institutional
architectures. The full set of data sources, variable constructions, and
supplementary results is deferred to Appendix~\ref{app:more_empirical_facts}.

\subsection{Symbolic Policy and Narrative Mobilization}
\label{subsec:d_narrative}

India's post-independence polity inherited a constitutional encouragement of
cow protection (Article~48) but no national prohibition. Through the first
four post-independence decades, beef regulation remained a state-level matter,
with Congress-governed states tolerating buffalo and low-yield cattle
slaughter and a handful of Jan Sangh- or BJP-governed states pushing stricter
enforcement. The legal landscape changed discontinuously after the BJP's
national victory in 2014. Between 2015 and 2017, a wave of state-level
legislation extended the ban from cows to the entire bovine family, reversed
the burden of proof, and raised maximum penalties to up to ten years'
imprisonment.\footnote{Maharashtra Animal Preservation (Amendment) Act, 2015;
Haryana Gauvansh Sanrakshan Act, 2015; analogous amendments in Gujarat,
Jharkhand, Uttar Pradesh, Madhya Pradesh, and Uttarakhand. A 2017 Ministry of
Environment notification sought to restrict cattle sales nationally before
being stayed by the Supreme Court.} \autoref{fig:ban_adoption} records the
cumulative adoption of six ban types across states, and Map~A
(\autoref{fig:maps}) summarizes cumulative ban intensity as of 2020.

\textbf{Majority-side narrative supply.} To trace the co-movement of the
symbolic instrument $d$ with majority-side mobilization $E_M$, I assemble a
corpus of $2{,}324$ BJP leadership texts from three sources: the party's
official speech archive
(\href{https://www.bjp.org/speeches}{bjp.org/speeches}, 1996--2025), the Prime
Minister's office
(\href{https://www.pmindia.gov.in}{pmindia.gov.in}, 2023--2026), and BJP
Members' floor speeches in the Lok Sabha, harvested from the official DSpace
backend of \href{https://elibrary.sansad.in}{elibrary.sansad.in}
(1996--2024). The Lok Sabha supplement fills the pre-2014 gap in the party's
online archive and captures parliamentary, as distinct from campaign,
register. Speakers include the Prime Minister, the four BJP National
Presidents of the period (Amit Shah, Rajnath Singh, J.~P.~Nadda, and Narendra
Modi himself), and senior front-benchers (Arun Jaitley, Sushma Swaraj, Nitin
Gadkari). Each speech is classified into one of four categories
by a keyword-lexicon procedure: (i)~\emph{identity-coded} content
(cow/beef/gau, Hindu--Muslim pairs, love jihad, CAA/NRC, ghar wapsi, Ayodhya,
UCC, hijab, loudspeaker, Gyanvapi, waqf, triple talaq, and cognate terms in
Hindi); (ii)~\emph{economic-programmatic} content (MSME, PM-KISAN, Ujjwala,
Aadhaar-linked DBT, GST, ``Make in India,'' ``Viksit Bharat,'' and cognates);
(iii)~\emph{security-and-foreign} content (Pakistan, Galwan, Article~370,
Operation Sindoor); (iv)~\emph{residual}.

\textbf{Minority-side narrative.} A symmetric majority-side corpus proved
infeasible: of the three major umbrella organizations --- the All India
Majlis-e-Ittehadul Muslimeen (AIMIM), the Jamiat Ulama-i-Hind, and the All
India Muslim Personal Law Board --- none operates a crawlable online speech
archive comparable to \texttt{bjp.org}.\footnote{As of April~2026,
\texttt{aimim.in} fails to resolve; \texttt{jamiatulama.org} resolves to a
parked GoDaddy lander; \texttt{aimplboard.in} retains a static early-2000s
site whose press-release endpoint returns~500. The archives of Owaisi's public
speeches persist on social-media platforms but are not systematically
catalogued with dates and transcripts. This asymmetry in narrative
infrastructure is itself informative: it is the observable expression of the
organizational-capacity parameter $\kappa$ in the Minority's collective-action
problem (\S\ref{subsec:joint_contest}), and tightens rather than weakens the
model's claim that the Minority's mobilization is supply-constrained.} I
therefore construct a minority-side salience series from a proxy corpus of
$3{,}783$ dated Indian news items, $2010$--$2025$, curated from major
English-language outlets --- The Hindu, The Times of India, NDTV, India Today,
Indian Express, Al Jazeera, Scroll, Maktoob, SabrangIndia, Muslim Mirror, and
the BBC --- using $65$ keyword queries covering AIMIM/Owaisi, Jamiat, AIMPLB,
cow-related lynchings (Akhlaq, Pehlu Khan, Junaid, Tabrez), Shaheen Bagh,
CAA/NRC, hijab, waqf, UCC, Gyanvapi, Babri, Bilkis, and Kashmir Muslim
leadership. Each item is classified into identity-coded, counter-mobilization
(protest, bandh, satyagraha, rally, ``not in my name,'' boycott),
security-and-foreign, economic, or residual.

\textbf{Results.} \autoref{fig:narrative_timeseries} plots the monthly share
of identity-coded content on the majority side (top panel) and the combined
identity-coded and counter-mobilization share on the minority side (bottom
panel), smoothed over a three-month rolling window. Four patterns are
visible. First, the majority-side identity share steps up in 2015--17 around
the state-level ban wave, contracts during the Covid period, and re-expands
in 2022--24 around the Uniform Civil Code and Waqf debates. Second, the
minority-side share is quiescent in 2010--14, rises discretely with the 2015
Dadri lynching and the 2017 Pehlu Khan killing, peaks during the 2019--20
CAA/NRC and Shaheen Bagh episode, and remains elevated through 2025. Third,
the two series move together at high frequency around legislative shocks
(2015 state bans, 2019 CAA), consistent with the joint contest structure of
Lemma~\ref{lem:contest_equilibrium} in which $d>0$ activates both
$E_M^{*}$~and~$E_m^{*}$. Fourth, the episode-level evidence is consistent
with the inverted-U bound on the Security Premium
(Lemma~\ref{lem:reaction_invertedU}): the majority-side identity share peaks
in 2016--17 and again in 2024, with intermediate compressions that coincide
with electorally costly overreach events (the 2018 Supreme Court stay of the
livestock-market notification; the 2020--21 farm laws, which were ultimately
repealed).

% ============================================================================
\begin{figure}[t]
    \centering
    \includegraphics[width=0.92\textwidth]{figures/empirics/narrative_timeseries.pdf}
    \caption{\normalsize\bf Majority and Minority Identity-Coded Narrative
    Share, India 2010--2025.}
    \floatfoot*{\footnotesize\textit{Notes.} Top panel: share of BJP
    leadership texts classified as identity-coded in the month,
    $N=2{,}324$ from \texttt{bjp.org/speeches}, \texttt{pmindia.gov.in}, and
    BJP Lok Sabha floor speeches harvested from
    \texttt{elibrary.sansad.in}, 1996m9--2026m4. Bottom panel: share of
    curated Muslim-voice items classified as identity-coded or
    counter-mobilization in the month, $N=3{,}792$ items
    ($3{,}783$~news-outlet proxies plus~$9$ AIMIM Lok~Sabha floor
    speeches), 2004m12--2025m12. Both series are three-month rolling
    means. Vertical lines mark the Dadri lynching (Sep~2015), the Haryana
    Gauvansh Sanrakshan Act (Mar~2015), the Pehlu Khan lynching (Apr~2017),
    the Citizenship (Amendment) Act passage and onset of Shaheen Bagh
    (Dec~2019), the farm laws (Sep~2020) and their repeal (Nov~2021), the
    Mahsa Amini episode in Iran (Sep~2022), and the 2024 Kanwar Yatra. The
    classification lexicon and methodology are reported in
    Appendix~\ref{app:narrative_classification}. Sources: bjp.org/speeches;
    pmindia.gov.in/en/news\_updates; curated news-outlet RSS archives.}
    \label{fig:narrative_timeseries}
\end{figure}
% ============================================================================

% ============================================================================
\begin{figure}[t]
    \centering
    \includegraphics[width=0.82\textwidth]{figures/empirics/map_ban_intensity.pdf}
    \caption{\normalsize\bf Composite Cattle-Slaughter Ban Intensity by
    State, Cumulative Through~2020.}
    \floatfoot*{\footnotesize\textit{Notes.} Composite ban intensity is
    computed as the sum over six ban types (cow, bullock, bull, buffalo
    slaughter; beef possession; beef sales) of
    $\max(\text{current year}-\text{adoption year},0)$, averaged over
    1990--2020. Higher values indicate earlier and broader adoption. Shapefile:
    GADM~v4.1 administrative level~1. Underlying policy
    data: Cattle Slaughter Policy dataset, constructed from state-level
    legislative records; see Appendix~\ref{app:ban_intensity}.}
    \label{fig:ban_adoption}
\end{figure}
% ============================================================================

\subsection{Cycles and Persistence in India's Political Competition}
\label{subsec:cycles}

The dynamic extension of the model in Section~\ref{sec:dynamic} delivers two
qualitatively distinct political paths: with sufficiently patient voters,
the Hawk becomes an absorbing state of the polity (Proposition~\ref{prop:trap});
with myopic retrospective voters, the same mechanism generates recurrent
Hawk--Dove regime alternation (Proposition~\ref{prop:cycles}). India's
federal structure provides 30 quasi-independent state polities observed over
1984--2024, in which both paths can be distinguished simultaneously.

I construct a state-year panel of Vidhan Sabha (state legislative assembly)
outcomes by combining the Bhavnani India Election Dataset v2.0 (1977--2015)
with supplementary state-year records I harvested from the Election
Commission of India and Wikipedia for post-2015 elections (the $51$ Vidhan
Sabha elections and the 2019 and 2024 Lok Sabha elections held in the
extension window). The panel is forward-filled between elections, so each
state-year is assigned the incumbent party of the preceding Assembly
result. Using the BJP's seat share, the top-1 party's vote share margin, and
the incumbent party identity, I classify each state-year as
\emph{Hawk-dominant} if the BJP rules with either a seat-share majority or a
vote-share margin of at least~$10$~pp;
\emph{Dove-dominant} if a non-BJP party rules with a margin of at
least~$10$~pp; and \emph{competitive} otherwise.

\autoref{fig:cycles_four_panel} presents four panels derived from the panel.
Panel~(A) plots the national BJP vote and seat shares in Lok Sabha elections
from~1984 to~2024. Two step increases stand out: the 1999--2004 NDA years
and the post-2014 Modi wave, with a peak of~55.8\%~of seats in~2019 on a
$37.4\%$~vote share. Panel~(B) plots the fraction of state Assemblies
governed by the BJP across time: the share rises from near-zero before~1990
to~$0.55$ after~2017. Panel~(C) plots the distribution of uninterrupted
incumbency tenure at the state level, by party group. The BJP distribution
has a pronounced right tail: the party holds five state spells of at
least~11~consecutive years, including Gujarat 1995--2024 (30~years),
Madhya~Pradesh 2003--2018 (15~years), Chhattisgarh 2003--2018 (15~years), and
Uttar~Pradesh 1991--2001 (11~years). The INC and regional-party
distributions are centred on shorter tenures. Panel~(D) plots the state-year
regime heatmap. Three patterns are visible. First, the post-2014 Hindi Belt
--- U.P., M.P., Rajasthan, Chhattisgarh, Gujarat, Haryana, Jharkhand ---
is populated by long Hawk-dominant spells with few transitions, the
empirical footprint of the conflict trap. Second, the southern and eastern
states --- Tamil~Nadu, Kerala, West~Bengal, Andhra~Pradesh, Telangana ---
remain Dove-dominant throughout, consistent with $\delta>\bar{\delta}$ not
applying uniformly. Third, an intermediate set of states --- Himachal~Pradesh
(seven regime transitions), Rajasthan ($6$), Haryana ($5$), Karnataka~($4$),
and Delhi --- alternate between categories, the empirical footprint of the
cycles case.

% ============================================================================
\begin{figure}[t]
    \centering
    \includegraphics[width=0.98\textwidth]{figures/empirics/cycles_four_panel.pdf}
    \caption{\normalsize\bf Cycles and Persistence in Indian Electoral
    Competition, 1984--2024.}
    \floatfoot*{\footnotesize\textit{Notes.} (A)~National BJP vote share and
    Lok Sabha seat share, 1984--2024. (B)~Fraction of 30~Indian states (and
    NCT of Delhi) governed by the BJP or a BJP-led coalition in each year.
    (C)~Histogram of uninterrupted Assembly-level incumbency tenure by party
    group \{BJP, INC, regional/other\}; spells that are still ongoing at the
    end of the panel are right-truncated. (D)~State-year regime classification
    (rows: $30$ states, alphabetical; columns: 1984--2024). A state-year is
    \emph{Hawk-dominant} when the BJP rules with seat-share
    $\geq 50\%$~or~vote-share margin~$\geq 10$~pp;
    \emph{Dove-dominant} when a non-BJP party rules with margin
    $\geq 10$~pp; \emph{competitive} otherwise.
    Data: Bhavnani India Election Dataset v2.0 (1977--2015) and
    author-compiled supplementary panel from Wikipedia and ECI for the
    $51$~Vidhan Sabha elections and 2019, 2024 Lok Sabha elections
    held in 2015--2025. See Appendix~\ref{app:data_sources} for the
    construction and sanity checks.}
    \label{fig:cycles_four_panel}
\end{figure}
% ============================================================================

The co-existence of long Hawk spells and intermediate-state cycles within the
same federation is not, by itself, a test of
Propositions~\ref{prop:trap} or~\ref{prop:cycles}: the classification is
reduced-form, and cross-state heterogeneity in~$\delta$, in the moderate-
defection technology $\chi_{\mathrm{mod}}(d)$, and in the primitive
state-capacity gap $S_A/S_B$ will all push states toward different paths.
What the panel does rule out is a uniform story --- either global
entrenchment or global alternation --- and is consistent with the model's
joint prediction that the same mechanism can produce both outcomes along
different parameter paths. Gujarat's 30-year Hawk spell, in particular, is the
archetype of the conflict trap: a continuous stretch in which the BJP
retained power through the 2002 communal violence and its aftermath, three
Lok Sabha electoral cycles, and at least one episode of sharp economic
contraction (the 2002--03 slowdown). That it did so while retaining
Dove-dominant neighbours (Rajasthan pre-2013, the 2018 swing) is exactly the
cross-state heterogeneity the model predicts.

\subsection{Voter Welfare and Political Dominance}
\label{subsec:voter_space}

Proposition~\ref{prop:provocation} states that the Hawk's electoral survival
requires turning the joint contest on when~$\Delta u>0$, and that the
symbolic-policy path is chosen to substitute an activated Security Premium
for the Hawk's economic deficit. The empirical footprint of this
substitution is a state-year distribution in the $(\Delta u,\text{BJP
dominance})$~space in which states with high cumulative ban intensity cluster
in a region distinct from states that have not used symbolic policy. This
subsection marks that distribution.

\textbf{Measurement.} I proxy $\Delta u$ by state GDP per capita relative to
the national average (where $1.0$ denotes the national average in the given
year), taken from the Reserve Bank of India's relative per-capita income
series at $1990, 2000, 2010, 2020$, and~$2023$ anchor years. I proxy BJP
dominance by the BJP's Vidhan Sabha seat share in the state-year, computed
from the election panel of \S\ref{subsec:cycles}. Composite ban intensity is
the state-level construction of \S\ref{subsec:d_narrative} carried forward
through~2020.

\textbf{Maps.} Maps~B and~C in \autoref{fig:maps} visualize the cross-section.
Map~A (ban intensity, already shown as \autoref{fig:ban_adoption}) clusters
on the northern and western states: Gujarat, Haryana, Punjab, Himachal
Pradesh, Uttarakhand, and large parts of Uttar Pradesh and Bihar. Map~B
plots Hindu--Muslim violence intensity per capita, pooling incidents from
the Varshney--Wilkinson dataset ($1950$--$1995$) and the GDELT Hindu--Muslim
variant ($1980$--$2024$), normalized by state population.\footnote{State
population is interpolated annually from Census~2011 shares rescaled by the
World~Bank national-population series; see
Appendix~\ref{app:population}.} The violence map co-locates with the ban map
along the Gangetic-plain belt and Gujarat. Map~C plots the change in BJP
state-Assembly vote share between the 2009-anchor and the 2019-anchor Vidhan
Sabha elections. The post-2014 breakthrough states ---Tripura (+42~pp),
Manipur (+35~pp), West~Bengal (+34~pp), Haryana (+27~pp), Uttar~Pradesh
(+23~pp), Assam (+22~pp), Odisha (+17~pp) --- are precisely the states that
had not previously delivered a BJP government and in which the post-2014 wave
represents an entry rather than a consolidation.

\textbf{Two-dimensional scatter.}
\autoref{fig:voter_space_2d} places each state-year in the $(\Delta
u,\text{BJP seat share})$ space, with ban intensity mapped to point size and
year to point color, for the anchor years $1990, 2000, 2010, 2020, 2023$.
Two centroids are marked. For the high-ban-intensity states (top tercile),
the $2000$~centroid sits at $(1.21, 0.27)$ --- above average relative income
and appreciable BJP dominance --- and the $2020$~centroid sits at
$(0.79, 0.33)$: relative income has fallen below the national average
while BJP dominance has risen. For the low-ban-intensity states (bottom
tercile), the $2000$~centroid sits at $(0.77, 0.01)$ and the $2020$~centroid
at $(0.41, 0.02)$: relative income has fallen sharply, while BJP dominance
is negligible throughout. The two arrows together are the empirical
signature of Proposition~\ref{prop:provocation}: in states where symbolic
policy was used intensively, political dominance stayed high (rose, in
fact) even as the relative-income position deteriorated; in states where
symbolic policy was not used, the collapse in relative income was not
offset. The relative movement in relative income is larger than the
movement in BJP dominance, consistent with the provocation-ceiling bound of
equation~\eqref{eq:delta_u_bar_bound}: the electoral offset provided by
activated symbolic contest is structurally limited, and even intense use of
$d$ cannot fully compensate for large income deficits.

% ============================================================================
\begin{figure}[t]
    \centering
    \subfigure[Composite cattle-slaughter ban intensity (cumulative through~2020).]{
        \includegraphics[width=0.32\textwidth]{figures/empirics/map_ban_intensity.pdf}
        \label{fig:map_ban_intensity}
    }
    \subfigure[Hindu--Muslim violence intensity per capita (1950--2024).]{
        \includegraphics[width=0.32\textwidth]{figures/empirics/map_conflict_intensity.pdf}
        \label{fig:map_conflict_intensity}
    }
    \subfigure[$\Delta$~BJP Vidhan Sabha vote share, 2009~$\to$~2019 (pp).]{
        \includegraphics[width=0.32\textwidth]{figures/empirics/map_bjp_delta_2009_2019.pdf}
        \label{fig:map_bjp_delta}
    }
    \caption{\normalsize\bf Spatial Distribution of Symbolic-Policy Intensity,
    Hindu--Muslim Violence, and BJP Breakthrough.}
    \floatfoot*{\footnotesize\textit{Notes.} (A)~Composite ban intensity as
    defined in \autoref{fig:ban_adoption}. (B)~Pooled Hindu--Muslim violent
    incidents per capita, sources: Varshney--Wilkinson 1950--1995 and GDELT
    Hindu--Muslim variant 1980--2024; population panel at
    Appendix~\ref{app:population}. (C)~BJP state-Assembly vote-share change
    between the Vidhan Sabha election closest to 2009 and the Vidhan Sabha
    election closest to 2019 for each state, in percentage points;
    diverging colour scale centred at zero.}
    \label{fig:maps}
\end{figure}
% ============================================================================

% ============================================================================
\begin{figure}[t]
    \centering
    \includegraphics[width=0.88\textwidth]{figures/empirics/voter_space_2d.pdf}
    \caption{\normalsize\bf Voter Welfare and Political Dominance, Indian
    States~1990--2023.}
    \floatfoot*{\footnotesize\textit{Notes.} Each point is a state-year
    observation. Horizontal axis: state GDP per capita relative to the
    national average (1.0~=~national average in the year); vertical axis:
    BJP state-Assembly seat share; point size: composite cumulative ban
    intensity (see \autoref{fig:ban_adoption}); point colour: year. Anchor
    years $\{1990, 2000, 2010, 2020, 2023\}$. Two arrows track the centroid
    movement of the top-tercile and bottom-tercile ban-intensity states
    between 2000 and 2020: high-ban states moved from $(1.21, 0.27)$ to
    $(0.79, 0.33)$; low-ban states moved from $(0.77, 0.01)$ to
    $(0.41, 0.02)$. The five states with the largest post-2014
    seat-share movement are labelled.}
    \label{fig:voter_space_2d}
\end{figure}
% ============================================================================

\subsection{Other Electoral Democracies: US, Colombia, and Iran}
\label{subsec:other_democracies}

The model is presented in the context of India's cow-slaughter and citizenship
legislation, but its primitives are generic to two-party electoral environments
with a salient out-group. Three further cases illustrate the mechanism
under different institutional architectures.

\emph{The United States, post-2015.}
The Republican Party's post-2015 turn to immigration enforcement and
culturally coded domestic security fits the model closely. The Democratic
comparative advantage lies in universalist economic provision (the Affordable
Care Act, the American Rescue Plan, the Inflation Reduction Act), consistent
with the Dove's role in $(S_B,\bar g_B)$. The symbolic instruments have been
legible: Executive Order~13769, border-wall construction, the 2018
family-separation policy, and the federalisation of culture-war education
policy. Each is a $d$-type policy with no direct $g$ payoff. The 2018
midterms, held immediately after the family-separation episode, returned a
Democratic House majority, consistent with
Lemma~\ref{lem:reaction_invertedU}'s inverted-U prediction: beyond
$d^{\mathrm{peak}}$, the Hawk's moderate-defection cost $\chi_{\mathrm{mod}}$
accumulates faster than the Security Premium. The dynamic path since~2016
---alternation~2016/2020/2024 --- is closer to the cycles case of
Proposition~\ref{prop:cycles} than to the trap of
Proposition~\ref{prop:trap}, consistent with shorter voter horizons and with
the institutional absence of the candidacy-gatekeeping that supports the
Indian and Iranian trap paths.

\emph{Colombia, the 2016 peace plebiscite.}
Colombia's post-2002 Uribista period extended the mechanism to an environment
in which the Minority is an ideologically coded armed party (the FARC) rather
than an ethnic or religious domestic out-group. The \emph{Democratic
Security} doctrine of Álvaro Uribe (2002--2010) and its continuation in the
Centro Democrático after~2014 fit the Hawk type: it promised visible coercive
capacity against FARC, imposed fiscal costs (the ``democratic security
tax''), and activated a two-sided contest in which paramilitary mobilisation
rose alongside state enforcement. The 2016 plebiscite on the Santos--FARC
peace accord is the sharpest case of Lemma~\ref{lem:dove_corner} cutting in
reverse: when the Dove attempted to move the system to $d=0$, enough voters
preferred the Hawk's manufactured-threat environment to block the transition.
That the successor administration subsequently implemented substantial
portions of the same accord demonstrates that the narrow plebiscite margin
was the binding electoral constraint, not a technical defect in the agreement.

\emph{Iran, morality policing and the 2022 Amini episode.}
The Islamic Republic of Iran adapts the model to a hybrid regime in which
competitive elections occur within candidacy-gatekeeping institutions (the
Guardian Council). The hard-line bloc associated with the Islamic
Revolutionary Guard Corps occupies the Hawk type; its most visible symbolic
instrument is the regulation of women's public conduct --- compulsory hijab,
the \emph{gasht-e ershad} morality police, and periodic ``bad hijab''
campaigns. The Mahsa Amini episode of September~2022 and the ensuing
\emph{Woman, Life, Freedom} protests are a textbook instance of the
joint-contest structure of Lemma~\ref{lem:contest_equilibrium}: a $d$~shock
triggered a large $E_m^*$, which in turn activated the regime's $S$~apparatus
(the Basij). The re-intensification of morality-police enforcement after the
protest's decline is the trap of Proposition~\ref{prop:trap} re-closing ---
but with a specifically institutional interpretation: the ``patience'' that
sustains the Hawk in Iran is supplied by candidacy gatekeeping rather than
by voter preferences over permanence, as the model's reading of
candidacy-filtered regimes anticipates.

Together, the three cases show that the India-centric analysis of
\S\S\ref{subsec:d_narrative}--\ref{subsec:voter_space} is an instance of a
more general pattern rather than a country-specific artefact. In each case
the incumbent's symbolic-policy instrument substitutes an activated contest
environment for a structural economic disadvantage, the inverted-U bound
disciplines the use of the instrument, and the dynamic path is shaped by the
interaction of voter patience, institutional gatekeeping, and the strength
of the moderate-defection channel.

% End of \section{Empirical Facts}
 between \section{Introduction} and
% \section{The Model}, i.e. around line 268 of model.tex.
% =============================================================================

\section{Empirical Facts} \label{sec:empirical_facts}

This section presents four empirical regularities that motivate the model and
anchor its primitives in the Indian case. Subsection~\ref{subsec:d_narrative}
documents the staggered intensification of the symbolic-policy instrument~$d$
and the co-movement of majority and minority narrative mobilization,
$E_M$~and~$E_m$. Subsection~\ref{subsec:cycles} shows the co-existence of
persistent Hawk-state absorption and recurrent regime cycles across Indian
states, the two dynamic paths identified by Propositions~\ref{prop:trap}
and~\ref{prop:cycles}. Subsection~\ref{subsec:voter_space} places states in the
two-dimensional space $(\Delta u,\text{BJP dominance})$ with ban intensity as a
third margin, tracing the empirical footprint of the Security Premium of
Lemma~\ref{lem:reaction_invertedU}. Subsection~\ref{subsec:other_democracies}
closes with three scope cases --- the United States, Colombia, and Iran ---
that illustrate how the same mechanism operates under different institutional
architectures. The full set of data sources, variable constructions, and
supplementary results is deferred to Appendix~\ref{app:more_empirical_facts}.

\subsection{Symbolic Policy and Narrative Mobilization}
\label{subsec:d_narrative}

India's post-independence polity inherited a constitutional encouragement of
cow protection (Article~48) but no national prohibition. Through the first
four post-independence decades, beef regulation remained a state-level matter,
with Congress-governed states tolerating buffalo and low-yield cattle
slaughter and a handful of Jan Sangh- or BJP-governed states pushing stricter
enforcement. The legal landscape changed discontinuously after the BJP's
national victory in 2014. Between 2015 and 2017, a wave of state-level
legislation extended the ban from cows to the entire bovine family, reversed
the burden of proof, and raised maximum penalties to up to ten years'
imprisonment.\footnote{Maharashtra Animal Preservation (Amendment) Act, 2015;
Haryana Gauvansh Sanrakshan Act, 2015; analogous amendments in Gujarat,
Jharkhand, Uttar Pradesh, Madhya Pradesh, and Uttarakhand. A 2017 Ministry of
Environment notification sought to restrict cattle sales nationally before
being stayed by the Supreme Court.} \autoref{fig:ban_adoption} records the
cumulative adoption of six ban types across states, and Map~A
(\autoref{fig:maps}) summarizes cumulative ban intensity as of 2020.

\textbf{Majority-side narrative supply.} To trace the co-movement of the
symbolic instrument $d$ with majority-side mobilization $E_M$, I assemble a
corpus of $2{,}324$ BJP leadership texts from three sources: the party's
official speech archive
(\href{https://www.bjp.org/speeches}{bjp.org/speeches}, 1996--2025), the Prime
Minister's office
(\href{https://www.pmindia.gov.in}{pmindia.gov.in}, 2023--2026), and BJP
Members' floor speeches in the Lok Sabha, harvested from the official DSpace
backend of \href{https://elibrary.sansad.in}{elibrary.sansad.in}
(1996--2024). The Lok Sabha supplement fills the pre-2014 gap in the party's
online archive and captures parliamentary, as distinct from campaign,
register. Speakers include the Prime Minister, the four BJP National
Presidents of the period (Amit Shah, Rajnath Singh, J.~P.~Nadda, and Narendra
Modi himself), and senior front-benchers (Arun Jaitley, Sushma Swaraj, Nitin
Gadkari). Each speech is classified into one of four categories
by a keyword-lexicon procedure: (i)~\emph{identity-coded} content
(cow/beef/gau, Hindu--Muslim pairs, love jihad, CAA/NRC, ghar wapsi, Ayodhya,
UCC, hijab, loudspeaker, Gyanvapi, waqf, triple talaq, and cognate terms in
Hindi); (ii)~\emph{economic-programmatic} content (MSME, PM-KISAN, Ujjwala,
Aadhaar-linked DBT, GST, ``Make in India,'' ``Viksit Bharat,'' and cognates);
(iii)~\emph{security-and-foreign} content (Pakistan, Galwan, Article~370,
Operation Sindoor); (iv)~\emph{residual}.

\textbf{Minority-side narrative.} A symmetric majority-side corpus proved
infeasible: of the three major umbrella organizations --- the All India
Majlis-e-Ittehadul Muslimeen (AIMIM), the Jamiat Ulama-i-Hind, and the All
India Muslim Personal Law Board --- none operates a crawlable online speech
archive comparable to \texttt{bjp.org}.\footnote{As of April~2026,
\texttt{aimim.in} fails to resolve; \texttt{jamiatulama.org} resolves to a
parked GoDaddy lander; \texttt{aimplboard.in} retains a static early-2000s
site whose press-release endpoint returns~500. The archives of Owaisi's public
speeches persist on social-media platforms but are not systematically
catalogued with dates and transcripts. This asymmetry in narrative
infrastructure is itself informative: it is the observable expression of the
organizational-capacity parameter $\kappa$ in the Minority's collective-action
problem (\S\ref{subsec:joint_contest}), and tightens rather than weakens the
model's claim that the Minority's mobilization is supply-constrained.} I
therefore construct a minority-side salience series from a proxy corpus of
$3{,}783$ dated Indian news items, $2010$--$2025$, curated from major
English-language outlets --- The Hindu, The Times of India, NDTV, India Today,
Indian Express, Al Jazeera, Scroll, Maktoob, SabrangIndia, Muslim Mirror, and
the BBC --- using $65$ keyword queries covering AIMIM/Owaisi, Jamiat, AIMPLB,
cow-related lynchings (Akhlaq, Pehlu Khan, Junaid, Tabrez), Shaheen Bagh,
CAA/NRC, hijab, waqf, UCC, Gyanvapi, Babri, Bilkis, and Kashmir Muslim
leadership. Each item is classified into identity-coded, counter-mobilization
(protest, bandh, satyagraha, rally, ``not in my name,'' boycott),
security-and-foreign, economic, or residual.

\textbf{Results.} \autoref{fig:narrative_timeseries} plots the monthly share
of identity-coded content on the majority side (top panel) and the combined
identity-coded and counter-mobilization share on the minority side (bottom
panel), smoothed over a three-month rolling window. Four patterns are
visible. First, the majority-side identity share steps up in 2015--17 around
the state-level ban wave, contracts during the Covid period, and re-expands
in 2022--24 around the Uniform Civil Code and Waqf debates. Second, the
minority-side share is quiescent in 2010--14, rises discretely with the 2015
Dadri lynching and the 2017 Pehlu Khan killing, peaks during the 2019--20
CAA/NRC and Shaheen Bagh episode, and remains elevated through 2025. Third,
the two series move together at high frequency around legislative shocks
(2015 state bans, 2019 CAA), consistent with the joint contest structure of
Lemma~\ref{lem:contest_equilibrium} in which $d>0$ activates both
$E_M^{*}$~and~$E_m^{*}$. Fourth, the episode-level evidence is consistent
with the inverted-U bound on the Security Premium
(Lemma~\ref{lem:reaction_invertedU}): the majority-side identity share peaks
in 2016--17 and again in 2024, with intermediate compressions that coincide
with electorally costly overreach events (the 2018 Supreme Court stay of the
livestock-market notification; the 2020--21 farm laws, which were ultimately
repealed).

% ============================================================================
\begin{figure}[t]
    \centering
    \includegraphics[width=0.92\textwidth]{figures/empirics/narrative_timeseries.pdf}
    \caption{\normalsize\bf Majority and Minority Identity-Coded Narrative
    Share, India 2010--2025.}
    \floatfoot*{\footnotesize\textit{Notes.} Top panel: share of BJP
    leadership texts classified as identity-coded in the month,
    $N=2{,}324$ from \texttt{bjp.org/speeches}, \texttt{pmindia.gov.in}, and
    BJP Lok Sabha floor speeches harvested from
    \texttt{elibrary.sansad.in}, 1996m9--2026m4. Bottom panel: share of
    curated Muslim-voice items classified as identity-coded or
    counter-mobilization in the month, $N=3{,}792$ items
    ($3{,}783$~news-outlet proxies plus~$9$ AIMIM Lok~Sabha floor
    speeches), 2004m12--2025m12. Both series are three-month rolling
    means. Vertical lines mark the Dadri lynching (Sep~2015), the Haryana
    Gauvansh Sanrakshan Act (Mar~2015), the Pehlu Khan lynching (Apr~2017),
    the Citizenship (Amendment) Act passage and onset of Shaheen Bagh
    (Dec~2019), the farm laws (Sep~2020) and their repeal (Nov~2021), the
    Mahsa Amini episode in Iran (Sep~2022), and the 2024 Kanwar Yatra. The
    classification lexicon and methodology are reported in
    Appendix~\ref{app:narrative_classification}. Sources: bjp.org/speeches;
    pmindia.gov.in/en/news\_updates; curated news-outlet RSS archives.}
    \label{fig:narrative_timeseries}
\end{figure}
% ============================================================================

% ============================================================================
\begin{figure}[t]
    \centering
    \includegraphics[width=0.82\textwidth]{figures/empirics/map_ban_intensity.pdf}
    \caption{\normalsize\bf Composite Cattle-Slaughter Ban Intensity by
    State, Cumulative Through~2020.}
    \floatfoot*{\footnotesize\textit{Notes.} Composite ban intensity is
    computed as the sum over six ban types (cow, bullock, bull, buffalo
    slaughter; beef possession; beef sales) of
    $\max(\text{current year}-\text{adoption year},0)$, averaged over
    1990--2020. Higher values indicate earlier and broader adoption. Shapefile:
    GADM~v4.1 administrative level~1. Underlying policy
    data: Cattle Slaughter Policy dataset, constructed from state-level
    legislative records; see Appendix~\ref{app:ban_intensity}.}
    \label{fig:ban_adoption}
\end{figure}
% ============================================================================

\subsection{Cycles and Persistence in India's Political Competition}
\label{subsec:cycles}

The dynamic extension of the model in Section~\ref{sec:dynamic} delivers two
qualitatively distinct political paths: with sufficiently patient voters,
the Hawk becomes an absorbing state of the polity (Proposition~\ref{prop:trap});
with myopic retrospective voters, the same mechanism generates recurrent
Hawk--Dove regime alternation (Proposition~\ref{prop:cycles}). India's
federal structure provides 30 quasi-independent state polities observed over
1984--2024, in which both paths can be distinguished simultaneously.

I construct a state-year panel of Vidhan Sabha (state legislative assembly)
outcomes by combining the Bhavnani India Election Dataset v2.0 (1977--2015)
with supplementary state-year records I harvested from the Election
Commission of India and Wikipedia for post-2015 elections (the $51$ Vidhan
Sabha elections and the 2019 and 2024 Lok Sabha elections held in the
extension window). The panel is forward-filled between elections, so each
state-year is assigned the incumbent party of the preceding Assembly
result. Using the BJP's seat share, the top-1 party's vote share margin, and
the incumbent party identity, I classify each state-year as
\emph{Hawk-dominant} if the BJP rules with either a seat-share majority or a
vote-share margin of at least~$10$~pp;
\emph{Dove-dominant} if a non-BJP party rules with a margin of at
least~$10$~pp; and \emph{competitive} otherwise.

\autoref{fig:cycles_four_panel} presents four panels derived from the panel.
Panel~(A) plots the national BJP vote and seat shares in Lok Sabha elections
from~1984 to~2024. Two step increases stand out: the 1999--2004 NDA years
and the post-2014 Modi wave, with a peak of~55.8\%~of seats in~2019 on a
$37.4\%$~vote share. Panel~(B) plots the fraction of state Assemblies
governed by the BJP across time: the share rises from near-zero before~1990
to~$0.55$ after~2017. Panel~(C) plots the distribution of uninterrupted
incumbency tenure at the state level, by party group. The BJP distribution
has a pronounced right tail: the party holds five state spells of at
least~11~consecutive years, including Gujarat 1995--2024 (30~years),
Madhya~Pradesh 2003--2018 (15~years), Chhattisgarh 2003--2018 (15~years), and
Uttar~Pradesh 1991--2001 (11~years). The INC and regional-party
distributions are centred on shorter tenures. Panel~(D) plots the state-year
regime heatmap. Three patterns are visible. First, the post-2014 Hindi Belt
--- U.P., M.P., Rajasthan, Chhattisgarh, Gujarat, Haryana, Jharkhand ---
is populated by long Hawk-dominant spells with few transitions, the
empirical footprint of the conflict trap. Second, the southern and eastern
states --- Tamil~Nadu, Kerala, West~Bengal, Andhra~Pradesh, Telangana ---
remain Dove-dominant throughout, consistent with $\delta>\bar{\delta}$ not
applying uniformly. Third, an intermediate set of states --- Himachal~Pradesh
(seven regime transitions), Rajasthan ($6$), Haryana ($5$), Karnataka~($4$),
and Delhi --- alternate between categories, the empirical footprint of the
cycles case.

% ============================================================================
\begin{figure}[t]
    \centering
    \includegraphics[width=0.98\textwidth]{figures/empirics/cycles_four_panel.pdf}
    \caption{\normalsize\bf Cycles and Persistence in Indian Electoral
    Competition, 1984--2024.}
    \floatfoot*{\footnotesize\textit{Notes.} (A)~National BJP vote share and
    Lok Sabha seat share, 1984--2024. (B)~Fraction of 30~Indian states (and
    NCT of Delhi) governed by the BJP or a BJP-led coalition in each year.
    (C)~Histogram of uninterrupted Assembly-level incumbency tenure by party
    group \{BJP, INC, regional/other\}; spells that are still ongoing at the
    end of the panel are right-truncated. (D)~State-year regime classification
    (rows: $30$ states, alphabetical; columns: 1984--2024). A state-year is
    \emph{Hawk-dominant} when the BJP rules with seat-share
    $\geq 50\%$~or~vote-share margin~$\geq 10$~pp;
    \emph{Dove-dominant} when a non-BJP party rules with margin
    $\geq 10$~pp; \emph{competitive} otherwise.
    Data: Bhavnani India Election Dataset v2.0 (1977--2015) and
    author-compiled supplementary panel from Wikipedia and ECI for the
    $51$~Vidhan Sabha elections and 2019, 2024 Lok Sabha elections
    held in 2015--2025. See Appendix~\ref{app:data_sources} for the
    construction and sanity checks.}
    \label{fig:cycles_four_panel}
\end{figure}
% ============================================================================

The co-existence of long Hawk spells and intermediate-state cycles within the
same federation is not, by itself, a test of
Propositions~\ref{prop:trap} or~\ref{prop:cycles}: the classification is
reduced-form, and cross-state heterogeneity in~$\delta$, in the moderate-
defection technology $\chi_{\mathrm{mod}}(d)$, and in the primitive
state-capacity gap $S_A/S_B$ will all push states toward different paths.
What the panel does rule out is a uniform story --- either global
entrenchment or global alternation --- and is consistent with the model's
joint prediction that the same mechanism can produce both outcomes along
different parameter paths. Gujarat's 30-year Hawk spell, in particular, is the
archetype of the conflict trap: a continuous stretch in which the BJP
retained power through the 2002 communal violence and its aftermath, three
Lok Sabha electoral cycles, and at least one episode of sharp economic
contraction (the 2002--03 slowdown). That it did so while retaining
Dove-dominant neighbours (Rajasthan pre-2013, the 2018 swing) is exactly the
cross-state heterogeneity the model predicts.

\subsection{Voter Welfare and Political Dominance}
\label{subsec:voter_space}

Proposition~\ref{prop:provocation} states that the Hawk's electoral survival
requires turning the joint contest on when~$\Delta u>0$, and that the
symbolic-policy path is chosen to substitute an activated Security Premium
for the Hawk's economic deficit. The empirical footprint of this
substitution is a state-year distribution in the $(\Delta u,\text{BJP
dominance})$~space in which states with high cumulative ban intensity cluster
in a region distinct from states that have not used symbolic policy. This
subsection marks that distribution.

\textbf{Measurement.} I proxy $\Delta u$ by state GDP per capita relative to
the national average (where $1.0$ denotes the national average in the given
year), taken from the Reserve Bank of India's relative per-capita income
series at $1990, 2000, 2010, 2020$, and~$2023$ anchor years. I proxy BJP
dominance by the BJP's Vidhan Sabha seat share in the state-year, computed
from the election panel of \S\ref{subsec:cycles}. Composite ban intensity is
the state-level construction of \S\ref{subsec:d_narrative} carried forward
through~2020.

\textbf{Maps.} Maps~B and~C in \autoref{fig:maps} visualize the cross-section.
Map~A (ban intensity, already shown as \autoref{fig:ban_adoption}) clusters
on the northern and western states: Gujarat, Haryana, Punjab, Himachal
Pradesh, Uttarakhand, and large parts of Uttar Pradesh and Bihar. Map~B
plots Hindu--Muslim violence intensity per capita, pooling incidents from
the Varshney--Wilkinson dataset ($1950$--$1995$) and the GDELT Hindu--Muslim
variant ($1980$--$2024$), normalized by state population.\footnote{State
population is interpolated annually from Census~2011 shares rescaled by the
World~Bank national-population series; see
Appendix~\ref{app:population}.} The violence map co-locates with the ban map
along the Gangetic-plain belt and Gujarat. Map~C plots the change in BJP
state-Assembly vote share between the 2009-anchor and the 2019-anchor Vidhan
Sabha elections. The post-2014 breakthrough states ---Tripura (+42~pp),
Manipur (+35~pp), West~Bengal (+34~pp), Haryana (+27~pp), Uttar~Pradesh
(+23~pp), Assam (+22~pp), Odisha (+17~pp) --- are precisely the states that
had not previously delivered a BJP government and in which the post-2014 wave
represents an entry rather than a consolidation.

\textbf{Two-dimensional scatter.}
\autoref{fig:voter_space_2d} places each state-year in the $(\Delta
u,\text{BJP seat share})$ space, with ban intensity mapped to point size and
year to point color, for the anchor years $1990, 2000, 2010, 2020, 2023$.
Two centroids are marked. For the high-ban-intensity states (top tercile),
the $2000$~centroid sits at $(1.21, 0.27)$ --- above average relative income
and appreciable BJP dominance --- and the $2020$~centroid sits at
$(0.79, 0.33)$: relative income has fallen below the national average
while BJP dominance has risen. For the low-ban-intensity states (bottom
tercile), the $2000$~centroid sits at $(0.77, 0.01)$ and the $2020$~centroid
at $(0.41, 0.02)$: relative income has fallen sharply, while BJP dominance
is negligible throughout. The two arrows together are the empirical
signature of Proposition~\ref{prop:provocation}: in states where symbolic
policy was used intensively, political dominance stayed high (rose, in
fact) even as the relative-income position deteriorated; in states where
symbolic policy was not used, the collapse in relative income was not
offset. The relative movement in relative income is larger than the
movement in BJP dominance, consistent with the provocation-ceiling bound of
equation~\eqref{eq:delta_u_bar_bound}: the electoral offset provided by
activated symbolic contest is structurally limited, and even intense use of
$d$ cannot fully compensate for large income deficits.

% ============================================================================
\begin{figure}[t]
    \centering
    \subfigure[Composite cattle-slaughter ban intensity (cumulative through~2020).]{
        \includegraphics[width=0.32\textwidth]{figures/empirics/map_ban_intensity.pdf}
        \label{fig:map_ban_intensity}
    }
    \subfigure[Hindu--Muslim violence intensity per capita (1950--2024).]{
        \includegraphics[width=0.32\textwidth]{figures/empirics/map_conflict_intensity.pdf}
        \label{fig:map_conflict_intensity}
    }
    \subfigure[$\Delta$~BJP Vidhan Sabha vote share, 2009~$\to$~2019 (pp).]{
        \includegraphics[width=0.32\textwidth]{figures/empirics/map_bjp_delta_2009_2019.pdf}
        \label{fig:map_bjp_delta}
    }
    \caption{\normalsize\bf Spatial Distribution of Symbolic-Policy Intensity,
    Hindu--Muslim Violence, and BJP Breakthrough.}
    \floatfoot*{\footnotesize\textit{Notes.} (A)~Composite ban intensity as
    defined in \autoref{fig:ban_adoption}. (B)~Pooled Hindu--Muslim violent
    incidents per capita, sources: Varshney--Wilkinson 1950--1995 and GDELT
    Hindu--Muslim variant 1980--2024; population panel at
    Appendix~\ref{app:population}. (C)~BJP state-Assembly vote-share change
    between the Vidhan Sabha election closest to 2009 and the Vidhan Sabha
    election closest to 2019 for each state, in percentage points;
    diverging colour scale centred at zero.}
    \label{fig:maps}
\end{figure}
% ============================================================================

% ============================================================================
\begin{figure}[t]
    \centering
    \includegraphics[width=0.88\textwidth]{figures/empirics/voter_space_2d.pdf}
    \caption{\normalsize\bf Voter Welfare and Political Dominance, Indian
    States~1990--2023.}
    \floatfoot*{\footnotesize\textit{Notes.} Each point is a state-year
    observation. Horizontal axis: state GDP per capita relative to the
    national average (1.0~=~national average in the year); vertical axis:
    BJP state-Assembly seat share; point size: composite cumulative ban
    intensity (see \autoref{fig:ban_adoption}); point colour: year. Anchor
    years $\{1990, 2000, 2010, 2020, 2023\}$. Two arrows track the centroid
    movement of the top-tercile and bottom-tercile ban-intensity states
    between 2000 and 2020: high-ban states moved from $(1.21, 0.27)$ to
    $(0.79, 0.33)$; low-ban states moved from $(0.77, 0.01)$ to
    $(0.41, 0.02)$. The five states with the largest post-2014
    seat-share movement are labelled.}
    \label{fig:voter_space_2d}
\end{figure}
% ============================================================================

\subsection{Other Electoral Democracies: US, Colombia, and Iran}
\label{subsec:other_democracies}

The model is presented in the context of India's cow-slaughter and citizenship
legislation, but its primitives are generic to two-party electoral environments
with a salient out-group. Three further cases illustrate the mechanism
under different institutional architectures.

\emph{The United States, post-2015.}
The Republican Party's post-2015 turn to immigration enforcement and
culturally coded domestic security fits the model closely. The Democratic
comparative advantage lies in universalist economic provision (the Affordable
Care Act, the American Rescue Plan, the Inflation Reduction Act), consistent
with the Dove's role in $(S_B,\bar g_B)$. The symbolic instruments have been
legible: Executive Order~13769, border-wall construction, the 2018
family-separation policy, and the federalisation of culture-war education
policy. Each is a $d$-type policy with no direct $g$ payoff. The 2018
midterms, held immediately after the family-separation episode, returned a
Democratic House majority, consistent with
Lemma~\ref{lem:reaction_invertedU}'s inverted-U prediction: beyond
$d^{\mathrm{peak}}$, the Hawk's moderate-defection cost $\chi_{\mathrm{mod}}$
accumulates faster than the Security Premium. The dynamic path since~2016
---alternation~2016/2020/2024 --- is closer to the cycles case of
Proposition~\ref{prop:cycles} than to the trap of
Proposition~\ref{prop:trap}, consistent with shorter voter horizons and with
the institutional absence of the candidacy-gatekeeping that supports the
Indian and Iranian trap paths.

\emph{Colombia, the 2016 peace plebiscite.}
Colombia's post-2002 Uribista period extended the mechanism to an environment
in which the Minority is an ideologically coded armed party (the FARC) rather
than an ethnic or religious domestic out-group. The \emph{Democratic
Security} doctrine of Álvaro Uribe (2002--2010) and its continuation in the
Centro Democrático after~2014 fit the Hawk type: it promised visible coercive
capacity against FARC, imposed fiscal costs (the ``democratic security
tax''), and activated a two-sided contest in which paramilitary mobilisation
rose alongside state enforcement. The 2016 plebiscite on the Santos--FARC
peace accord is the sharpest case of Lemma~\ref{lem:dove_corner} cutting in
reverse: when the Dove attempted to move the system to $d=0$, enough voters
preferred the Hawk's manufactured-threat environment to block the transition.
That the successor administration subsequently implemented substantial
portions of the same accord demonstrates that the narrow plebiscite margin
was the binding electoral constraint, not a technical defect in the agreement.

\emph{Iran, morality policing and the 2022 Amini episode.}
The Islamic Republic of Iran adapts the model to a hybrid regime in which
competitive elections occur within candidacy-gatekeeping institutions (the
Guardian Council). The hard-line bloc associated with the Islamic
Revolutionary Guard Corps occupies the Hawk type; its most visible symbolic
instrument is the regulation of women's public conduct --- compulsory hijab,
the \emph{gasht-e ershad} morality police, and periodic ``bad hijab''
campaigns. The Mahsa Amini episode of September~2022 and the ensuing
\emph{Woman, Life, Freedom} protests are a textbook instance of the
joint-contest structure of Lemma~\ref{lem:contest_equilibrium}: a $d$~shock
triggered a large $E_m^*$, which in turn activated the regime's $S$~apparatus
(the Basij). The re-intensification of morality-police enforcement after the
protest's decline is the trap of Proposition~\ref{prop:trap} re-closing ---
but with a specifically institutional interpretation: the ``patience'' that
sustains the Hawk in Iran is supplied by candidacy gatekeeping rather than
by voter preferences over permanence, as the model's reading of
candidacy-filtered regimes anticipates.

Together, the three cases show that the India-centric analysis of
\S\S\ref{subsec:d_narrative}--\ref{subsec:voter_space} is an instance of a
more general pattern rather than a country-specific artefact. In each case
the incumbent's symbolic-policy instrument substitutes an activated contest
environment for a structural economic disadvantage, the inverted-U bound
disciplines the use of the instrument, and the dynamic path is shaped by the
interaction of voter patience, institutional gatekeeping, and the strength
of the moderate-defection channel.

% End of \section{Empirical Facts}
 between \section{Introduction} and
% \section{The Model}, i.e. around line 268 of model.tex.
% =============================================================================

\section{Empirical Facts} \label{sec:empirical_facts}

This section presents four empirical regularities that motivate the model and
anchor its primitives in the Indian case. Subsection~\ref{subsec:d_narrative}
documents the staggered intensification of the symbolic-policy instrument~$d$
and the co-movement of majority and minority narrative mobilization,
$E_M$~and~$E_m$. Subsection~\ref{subsec:cycles} shows the co-existence of
persistent Hawk-state absorption and recurrent regime cycles across Indian
states, the two dynamic paths identified by Propositions~\ref{prop:trap}
and~\ref{prop:cycles}. Subsection~\ref{subsec:voter_space} places states in the
two-dimensional space $(\Delta u,\text{BJP dominance})$ with ban intensity as a
third margin, tracing the empirical footprint of the Security Premium of
Lemma~\ref{lem:reaction_invertedU}. Subsection~\ref{subsec:other_democracies}
closes with three scope cases --- the United States, Colombia, and Iran ---
that illustrate how the same mechanism operates under different institutional
architectures. The full set of data sources, variable constructions, and
supplementary results is deferred to Appendix~\ref{app:more_empirical_facts}.

\subsection{Symbolic Policy and Narrative Mobilization}
\label{subsec:d_narrative}

India's post-independence polity inherited a constitutional encouragement of
cow protection (Article~48) but no national prohibition. Through the first
four post-independence decades, beef regulation remained a state-level matter,
with Congress-governed states tolerating buffalo and low-yield cattle
slaughter and a handful of Jan Sangh- or BJP-governed states pushing stricter
enforcement. The legal landscape changed discontinuously after the BJP's
national victory in 2014. Between 2015 and 2017, a wave of state-level
legislation extended the ban from cows to the entire bovine family, reversed
the burden of proof, and raised maximum penalties to up to ten years'
imprisonment.\footnote{Maharashtra Animal Preservation (Amendment) Act, 2015;
Haryana Gauvansh Sanrakshan Act, 2015; analogous amendments in Gujarat,
Jharkhand, Uttar Pradesh, Madhya Pradesh, and Uttarakhand. A 2017 Ministry of
Environment notification sought to restrict cattle sales nationally before
being stayed by the Supreme Court.} \autoref{fig:ban_adoption} records the
cumulative adoption of six ban types across states, and Map~A
(\autoref{fig:maps}) summarizes cumulative ban intensity as of 2020.

\textbf{Majority-side narrative supply.} To trace the co-movement of the
symbolic instrument $d$ with majority-side mobilization $E_M$, I assemble a
corpus of $2{,}324$ BJP leadership texts from three sources: the party's
official speech archive
(\href{https://www.bjp.org/speeches}{bjp.org/speeches}, 1996--2025), the Prime
Minister's office
(\href{https://www.pmindia.gov.in}{pmindia.gov.in}, 2023--2026), and BJP
Members' floor speeches in the Lok Sabha, harvested from the official DSpace
backend of \href{https://elibrary.sansad.in}{elibrary.sansad.in}
(1996--2024). The Lok Sabha supplement fills the pre-2014 gap in the party's
online archive and captures parliamentary, as distinct from campaign,
register. Speakers include the Prime Minister, the four BJP National
Presidents of the period (Amit Shah, Rajnath Singh, J.~P.~Nadda, and Narendra
Modi himself), and senior front-benchers (Arun Jaitley, Sushma Swaraj, Nitin
Gadkari). Each speech is classified into one of four categories
by a keyword-lexicon procedure: (i)~\emph{identity-coded} content
(cow/beef/gau, Hindu--Muslim pairs, love jihad, CAA/NRC, ghar wapsi, Ayodhya,
UCC, hijab, loudspeaker, Gyanvapi, waqf, triple talaq, and cognate terms in
Hindi); (ii)~\emph{economic-programmatic} content (MSME, PM-KISAN, Ujjwala,
Aadhaar-linked DBT, GST, ``Make in India,'' ``Viksit Bharat,'' and cognates);
(iii)~\emph{security-and-foreign} content (Pakistan, Galwan, Article~370,
Operation Sindoor); (iv)~\emph{residual}.

\textbf{Minority-side narrative.} A symmetric majority-side corpus proved
infeasible: of the three major umbrella organizations --- the All India
Majlis-e-Ittehadul Muslimeen (AIMIM), the Jamiat Ulama-i-Hind, and the All
India Muslim Personal Law Board --- none operates a crawlable online speech
archive comparable to \texttt{bjp.org}.\footnote{As of April~2026,
\texttt{aimim.in} fails to resolve; \texttt{jamiatulama.org} resolves to a
parked GoDaddy lander; \texttt{aimplboard.in} retains a static early-2000s
site whose press-release endpoint returns~500. The archives of Owaisi's public
speeches persist on social-media platforms but are not systematically
catalogued with dates and transcripts. This asymmetry in narrative
infrastructure is itself informative: it is the observable expression of the
organizational-capacity parameter $\kappa$ in the Minority's collective-action
problem (\S\ref{subsec:joint_contest}), and tightens rather than weakens the
model's claim that the Minority's mobilization is supply-constrained.} I
therefore construct a minority-side salience series from a proxy corpus of
$3{,}783$ dated Indian news items, $2010$--$2025$, curated from major
English-language outlets --- The Hindu, The Times of India, NDTV, India Today,
Indian Express, Al Jazeera, Scroll, Maktoob, SabrangIndia, Muslim Mirror, and
the BBC --- using $65$ keyword queries covering AIMIM/Owaisi, Jamiat, AIMPLB,
cow-related lynchings (Akhlaq, Pehlu Khan, Junaid, Tabrez), Shaheen Bagh,
CAA/NRC, hijab, waqf, UCC, Gyanvapi, Babri, Bilkis, and Kashmir Muslim
leadership. Each item is classified into identity-coded, counter-mobilization
(protest, bandh, satyagraha, rally, ``not in my name,'' boycott),
security-and-foreign, economic, or residual.

\textbf{Results.} \autoref{fig:narrative_timeseries} plots the monthly share
of identity-coded content on the majority side (top panel) and the combined
identity-coded and counter-mobilization share on the minority side (bottom
panel), smoothed over a three-month rolling window. Four patterns are
visible. First, the majority-side identity share steps up in 2015--17 around
the state-level ban wave, contracts during the Covid period, and re-expands
in 2022--24 around the Uniform Civil Code and Waqf debates. Second, the
minority-side share is quiescent in 2010--14, rises discretely with the 2015
Dadri lynching and the 2017 Pehlu Khan killing, peaks during the 2019--20
CAA/NRC and Shaheen Bagh episode, and remains elevated through 2025. Third,
the two series move together at high frequency around legislative shocks
(2015 state bans, 2019 CAA), consistent with the joint contest structure of
Lemma~\ref{lem:contest_equilibrium} in which $d>0$ activates both
$E_M^{*}$~and~$E_m^{*}$. Fourth, the episode-level evidence is consistent
with the inverted-U bound on the Security Premium
(Lemma~\ref{lem:reaction_invertedU}): the majority-side identity share peaks
in 2016--17 and again in 2024, with intermediate compressions that coincide
with electorally costly overreach events (the 2018 Supreme Court stay of the
livestock-market notification; the 2020--21 farm laws, which were ultimately
repealed).

% ============================================================================
\begin{figure}[t]
    \centering
    \includegraphics[width=0.92\textwidth]{figures/empirics/narrative_timeseries.pdf}
    \caption{\normalsize\bf Majority and Minority Identity-Coded Narrative
    Share, India 2010--2025.}
    \floatfoot*{\footnotesize\textit{Notes.} Top panel: share of BJP
    leadership texts classified as identity-coded in the month,
    $N=2{,}324$ from \texttt{bjp.org/speeches}, \texttt{pmindia.gov.in}, and
    BJP Lok Sabha floor speeches harvested from
    \texttt{elibrary.sansad.in}, 1996m9--2026m4. Bottom panel: share of
    curated Muslim-voice items classified as identity-coded or
    counter-mobilization in the month, $N=3{,}792$ items
    ($3{,}783$~news-outlet proxies plus~$9$ AIMIM Lok~Sabha floor
    speeches), 2004m12--2025m12. Both series are three-month rolling
    means. Vertical lines mark the Dadri lynching (Sep~2015), the Haryana
    Gauvansh Sanrakshan Act (Mar~2015), the Pehlu Khan lynching (Apr~2017),
    the Citizenship (Amendment) Act passage and onset of Shaheen Bagh
    (Dec~2019), the farm laws (Sep~2020) and their repeal (Nov~2021), the
    Mahsa Amini episode in Iran (Sep~2022), and the 2024 Kanwar Yatra. The
    classification lexicon and methodology are reported in
    Appendix~\ref{app:narrative_classification}. Sources: bjp.org/speeches;
    pmindia.gov.in/en/news\_updates; curated news-outlet RSS archives.}
    \label{fig:narrative_timeseries}
\end{figure}
% ============================================================================

% ============================================================================
\begin{figure}[t]
    \centering
    \includegraphics[width=0.82\textwidth]{figures/empirics/map_ban_intensity.pdf}
    \caption{\normalsize\bf Composite Cattle-Slaughter Ban Intensity by
    State, Cumulative Through~2020.}
    \floatfoot*{\footnotesize\textit{Notes.} Composite ban intensity is
    computed as the sum over six ban types (cow, bullock, bull, buffalo
    slaughter; beef possession; beef sales) of
    $\max(\text{current year}-\text{adoption year},0)$, averaged over
    1990--2020. Higher values indicate earlier and broader adoption. Shapefile:
    GADM~v4.1 administrative level~1. Underlying policy
    data: Cattle Slaughter Policy dataset, constructed from state-level
    legislative records; see Appendix~\ref{app:ban_intensity}.}
    \label{fig:ban_adoption}
\end{figure}
% ============================================================================

\subsection{Cycles and Persistence in India's Political Competition}
\label{subsec:cycles}

The dynamic extension of the model in Section~\ref{sec:dynamic} delivers two
qualitatively distinct political paths: with sufficiently patient voters,
the Hawk becomes an absorbing state of the polity (Proposition~\ref{prop:trap});
with myopic retrospective voters, the same mechanism generates recurrent
Hawk--Dove regime alternation (Proposition~\ref{prop:cycles}). India's
federal structure provides 30 quasi-independent state polities observed over
1984--2024, in which both paths can be distinguished simultaneously.

I construct a state-year panel of Vidhan Sabha (state legislative assembly)
outcomes by combining the Bhavnani India Election Dataset v2.0 (1977--2015)
with supplementary state-year records I harvested from the Election
Commission of India and Wikipedia for post-2015 elections (the $51$ Vidhan
Sabha elections and the 2019 and 2024 Lok Sabha elections held in the
extension window). The panel is forward-filled between elections, so each
state-year is assigned the incumbent party of the preceding Assembly
result. Using the BJP's seat share, the top-1 party's vote share margin, and
the incumbent party identity, I classify each state-year as
\emph{Hawk-dominant} if the BJP rules with either a seat-share majority or a
vote-share margin of at least~$10$~pp;
\emph{Dove-dominant} if a non-BJP party rules with a margin of at
least~$10$~pp; and \emph{competitive} otherwise.

\autoref{fig:cycles_four_panel} presents four panels derived from the panel.
Panel~(A) plots the national BJP vote and seat shares in Lok Sabha elections
from~1984 to~2024. Two step increases stand out: the 1999--2004 NDA years
and the post-2014 Modi wave, with a peak of~55.8\%~of seats in~2019 on a
$37.4\%$~vote share. Panel~(B) plots the fraction of state Assemblies
governed by the BJP across time: the share rises from near-zero before~1990
to~$0.55$ after~2017. Panel~(C) plots the distribution of uninterrupted
incumbency tenure at the state level, by party group. The BJP distribution
has a pronounced right tail: the party holds five state spells of at
least~11~consecutive years, including Gujarat 1995--2024 (30~years),
Madhya~Pradesh 2003--2018 (15~years), Chhattisgarh 2003--2018 (15~years), and
Uttar~Pradesh 1991--2001 (11~years). The INC and regional-party
distributions are centred on shorter tenures. Panel~(D) plots the state-year
regime heatmap. Three patterns are visible. First, the post-2014 Hindi Belt
--- U.P., M.P., Rajasthan, Chhattisgarh, Gujarat, Haryana, Jharkhand ---
is populated by long Hawk-dominant spells with few transitions, the
empirical footprint of the conflict trap. Second, the southern and eastern
states --- Tamil~Nadu, Kerala, West~Bengal, Andhra~Pradesh, Telangana ---
remain Dove-dominant throughout, consistent with $\delta>\bar{\delta}$ not
applying uniformly. Third, an intermediate set of states --- Himachal~Pradesh
(seven regime transitions), Rajasthan ($6$), Haryana ($5$), Karnataka~($4$),
and Delhi --- alternate between categories, the empirical footprint of the
cycles case.

% ============================================================================
\begin{figure}[t]
    \centering
    \includegraphics[width=0.98\textwidth]{figures/empirics/cycles_four_panel.pdf}
    \caption{\normalsize\bf Cycles and Persistence in Indian Electoral
    Competition, 1984--2024.}
    \floatfoot*{\footnotesize\textit{Notes.} (A)~National BJP vote share and
    Lok Sabha seat share, 1984--2024. (B)~Fraction of 30~Indian states (and
    NCT of Delhi) governed by the BJP or a BJP-led coalition in each year.
    (C)~Histogram of uninterrupted Assembly-level incumbency tenure by party
    group \{BJP, INC, regional/other\}; spells that are still ongoing at the
    end of the panel are right-truncated. (D)~State-year regime classification
    (rows: $30$ states, alphabetical; columns: 1984--2024). A state-year is
    \emph{Hawk-dominant} when the BJP rules with seat-share
    $\geq 50\%$~or~vote-share margin~$\geq 10$~pp;
    \emph{Dove-dominant} when a non-BJP party rules with margin
    $\geq 10$~pp; \emph{competitive} otherwise.
    Data: Bhavnani India Election Dataset v2.0 (1977--2015) and
    author-compiled supplementary panel from Wikipedia and ECI for the
    $51$~Vidhan Sabha elections and 2019, 2024 Lok Sabha elections
    held in 2015--2025. See Appendix~\ref{app:data_sources} for the
    construction and sanity checks.}
    \label{fig:cycles_four_panel}
\end{figure}
% ============================================================================

The co-existence of long Hawk spells and intermediate-state cycles within the
same federation is not, by itself, a test of
Propositions~\ref{prop:trap} or~\ref{prop:cycles}: the classification is
reduced-form, and cross-state heterogeneity in~$\delta$, in the moderate-
defection technology $\chi_{\mathrm{mod}}(d)$, and in the primitive
state-capacity gap $S_A/S_B$ will all push states toward different paths.
What the panel does rule out is a uniform story --- either global
entrenchment or global alternation --- and is consistent with the model's
joint prediction that the same mechanism can produce both outcomes along
different parameter paths. Gujarat's 30-year Hawk spell, in particular, is the
archetype of the conflict trap: a continuous stretch in which the BJP
retained power through the 2002 communal violence and its aftermath, three
Lok Sabha electoral cycles, and at least one episode of sharp economic
contraction (the 2002--03 slowdown). That it did so while retaining
Dove-dominant neighbours (Rajasthan pre-2013, the 2018 swing) is exactly the
cross-state heterogeneity the model predicts.

\subsection{Voter Welfare and Political Dominance}
\label{subsec:voter_space}

Proposition~\ref{prop:provocation} states that the Hawk's electoral survival
requires turning the joint contest on when~$\Delta u>0$, and that the
symbolic-policy path is chosen to substitute an activated Security Premium
for the Hawk's economic deficit. The empirical footprint of this
substitution is a state-year distribution in the $(\Delta u,\text{BJP
dominance})$~space in which states with high cumulative ban intensity cluster
in a region distinct from states that have not used symbolic policy. This
subsection marks that distribution.

\textbf{Measurement.} I proxy $\Delta u$ by state GDP per capita relative to
the national average (where $1.0$ denotes the national average in the given
year), taken from the Reserve Bank of India's relative per-capita income
series at $1990, 2000, 2010, 2020$, and~$2023$ anchor years. I proxy BJP
dominance by the BJP's Vidhan Sabha seat share in the state-year, computed
from the election panel of \S\ref{subsec:cycles}. Composite ban intensity is
the state-level construction of \S\ref{subsec:d_narrative} carried forward
through~2020.

\textbf{Maps.} Maps~B and~C in \autoref{fig:maps} visualize the cross-section.
Map~A (ban intensity, already shown as \autoref{fig:ban_adoption}) clusters
on the northern and western states: Gujarat, Haryana, Punjab, Himachal
Pradesh, Uttarakhand, and large parts of Uttar Pradesh and Bihar. Map~B
plots Hindu--Muslim violence intensity per capita, pooling incidents from
the Varshney--Wilkinson dataset ($1950$--$1995$) and the GDELT Hindu--Muslim
variant ($1980$--$2024$), normalized by state population.\footnote{State
population is interpolated annually from Census~2011 shares rescaled by the
World~Bank national-population series; see
Appendix~\ref{app:population}.} The violence map co-locates with the ban map
along the Gangetic-plain belt and Gujarat. Map~C plots the change in BJP
state-Assembly vote share between the 2009-anchor and the 2019-anchor Vidhan
Sabha elections. The post-2014 breakthrough states ---Tripura (+42~pp),
Manipur (+35~pp), West~Bengal (+34~pp), Haryana (+27~pp), Uttar~Pradesh
(+23~pp), Assam (+22~pp), Odisha (+17~pp) --- are precisely the states that
had not previously delivered a BJP government and in which the post-2014 wave
represents an entry rather than a consolidation.

\textbf{Two-dimensional scatter.}
\autoref{fig:voter_space_2d} places each state-year in the $(\Delta
u,\text{BJP seat share})$ space, with ban intensity mapped to point size and
year to point color, for the anchor years $1990, 2000, 2010, 2020, 2023$.
Two centroids are marked. For the high-ban-intensity states (top tercile),
the $2000$~centroid sits at $(1.21, 0.27)$ --- above average relative income
and appreciable BJP dominance --- and the $2020$~centroid sits at
$(0.79, 0.33)$: relative income has fallen below the national average
while BJP dominance has risen. For the low-ban-intensity states (bottom
tercile), the $2000$~centroid sits at $(0.77, 0.01)$ and the $2020$~centroid
at $(0.41, 0.02)$: relative income has fallen sharply, while BJP dominance
is negligible throughout. The two arrows together are the empirical
signature of Proposition~\ref{prop:provocation}: in states where symbolic
policy was used intensively, political dominance stayed high (rose, in
fact) even as the relative-income position deteriorated; in states where
symbolic policy was not used, the collapse in relative income was not
offset. The relative movement in relative income is larger than the
movement in BJP dominance, consistent with the provocation-ceiling bound of
equation~\eqref{eq:delta_u_bar_bound}: the electoral offset provided by
activated symbolic contest is structurally limited, and even intense use of
$d$ cannot fully compensate for large income deficits.

% ============================================================================
\begin{figure}[t]
    \centering
    \subfigure[Composite cattle-slaughter ban intensity (cumulative through~2020).]{
        \includegraphics[width=0.32\textwidth]{figures/empirics/map_ban_intensity.pdf}
        \label{fig:map_ban_intensity}
    }
    \subfigure[Hindu--Muslim violence intensity per capita (1950--2024).]{
        \includegraphics[width=0.32\textwidth]{figures/empirics/map_conflict_intensity.pdf}
        \label{fig:map_conflict_intensity}
    }
    \subfigure[$\Delta$~BJP Vidhan Sabha vote share, 2009~$\to$~2019 (pp).]{
        \includegraphics[width=0.32\textwidth]{figures/empirics/map_bjp_delta_2009_2019.pdf}
        \label{fig:map_bjp_delta}
    }
    \caption{\normalsize\bf Spatial Distribution of Symbolic-Policy Intensity,
    Hindu--Muslim Violence, and BJP Breakthrough.}
    \floatfoot*{\footnotesize\textit{Notes.} (A)~Composite ban intensity as
    defined in \autoref{fig:ban_adoption}. (B)~Pooled Hindu--Muslim violent
    incidents per capita, sources: Varshney--Wilkinson 1950--1995 and GDELT
    Hindu--Muslim variant 1980--2024; population panel at
    Appendix~\ref{app:population}. (C)~BJP state-Assembly vote-share change
    between the Vidhan Sabha election closest to 2009 and the Vidhan Sabha
    election closest to 2019 for each state, in percentage points;
    diverging colour scale centred at zero.}
    \label{fig:maps}
\end{figure}
% ============================================================================

% ============================================================================
\begin{figure}[t]
    \centering
    \includegraphics[width=0.88\textwidth]{figures/empirics/voter_space_2d.pdf}
    \caption{\normalsize\bf Voter Welfare and Political Dominance, Indian
    States~1990--2023.}
    \floatfoot*{\footnotesize\textit{Notes.} Each point is a state-year
    observation. Horizontal axis: state GDP per capita relative to the
    national average (1.0~=~national average in the year); vertical axis:
    BJP state-Assembly seat share; point size: composite cumulative ban
    intensity (see \autoref{fig:ban_adoption}); point colour: year. Anchor
    years $\{1990, 2000, 2010, 2020, 2023\}$. Two arrows track the centroid
    movement of the top-tercile and bottom-tercile ban-intensity states
    between 2000 and 2020: high-ban states moved from $(1.21, 0.27)$ to
    $(0.79, 0.33)$; low-ban states moved from $(0.77, 0.01)$ to
    $(0.41, 0.02)$. The five states with the largest post-2014
    seat-share movement are labelled.}
    \label{fig:voter_space_2d}
\end{figure}
% ============================================================================

\subsection{Other Electoral Democracies: US, Colombia, and Iran}
\label{subsec:other_democracies}

The model is presented in the context of India's cow-slaughter and citizenship
legislation, but its primitives are generic to two-party electoral environments
with a salient out-group. Three further cases illustrate the mechanism
under different institutional architectures.

\emph{The United States, post-2015.}
The Republican Party's post-2015 turn to immigration enforcement and
culturally coded domestic security fits the model closely. The Democratic
comparative advantage lies in universalist economic provision (the Affordable
Care Act, the American Rescue Plan, the Inflation Reduction Act), consistent
with the Dove's role in $(S_B,\bar g_B)$. The symbolic instruments have been
legible: Executive Order~13769, border-wall construction, the 2018
family-separation policy, and the federalisation of culture-war education
policy. Each is a $d$-type policy with no direct $g$ payoff. The 2018
midterms, held immediately after the family-separation episode, returned a
Democratic House majority, consistent with
Lemma~\ref{lem:reaction_invertedU}'s inverted-U prediction: beyond
$d^{\mathrm{peak}}$, the Hawk's moderate-defection cost $\chi_{\mathrm{mod}}$
accumulates faster than the Security Premium. The dynamic path since~2016
---alternation~2016/2020/2024 --- is closer to the cycles case of
Proposition~\ref{prop:cycles} than to the trap of
Proposition~\ref{prop:trap}, consistent with shorter voter horizons and with
the institutional absence of the candidacy-gatekeeping that supports the
Indian and Iranian trap paths.

\emph{Colombia, the 2016 peace plebiscite.}
Colombia's post-2002 Uribista period extended the mechanism to an environment
in which the Minority is an ideologically coded armed party (the FARC) rather
than an ethnic or religious domestic out-group. The \emph{Democratic
Security} doctrine of Álvaro Uribe (2002--2010) and its continuation in the
Centro Democrático after~2014 fit the Hawk type: it promised visible coercive
capacity against FARC, imposed fiscal costs (the ``democratic security
tax''), and activated a two-sided contest in which paramilitary mobilisation
rose alongside state enforcement. The 2016 plebiscite on the Santos--FARC
peace accord is the sharpest case of Lemma~\ref{lem:dove_corner} cutting in
reverse: when the Dove attempted to move the system to $d=0$, enough voters
preferred the Hawk's manufactured-threat environment to block the transition.
That the successor administration subsequently implemented substantial
portions of the same accord demonstrates that the narrow plebiscite margin
was the binding electoral constraint, not a technical defect in the agreement.

\emph{Iran, morality policing and the 2022 Amini episode.}
The Islamic Republic of Iran adapts the model to a hybrid regime in which
competitive elections occur within candidacy-gatekeeping institutions (the
Guardian Council). The hard-line bloc associated with the Islamic
Revolutionary Guard Corps occupies the Hawk type; its most visible symbolic
instrument is the regulation of women's public conduct --- compulsory hijab,
the \emph{gasht-e ershad} morality police, and periodic ``bad hijab''
campaigns. The Mahsa Amini episode of September~2022 and the ensuing
\emph{Woman, Life, Freedom} protests are a textbook instance of the
joint-contest structure of Lemma~\ref{lem:contest_equilibrium}: a $d$~shock
triggered a large $E_m^*$, which in turn activated the regime's $S$~apparatus
(the Basij). The re-intensification of morality-police enforcement after the
protest's decline is the trap of Proposition~\ref{prop:trap} re-closing ---
but with a specifically institutional interpretation: the ``patience'' that
sustains the Hawk in Iran is supplied by candidacy gatekeeping rather than
by voter preferences over permanence, as the model's reading of
candidacy-filtered regimes anticipates.

Together, the three cases show that the India-centric analysis of
\S\S\ref{subsec:d_narrative}--\ref{subsec:voter_space} is an instance of a
more general pattern rather than a country-specific artefact. In each case
the incumbent's symbolic-policy instrument substitutes an activated contest
environment for a structural economic disadvantage, the inverted-U bound
disciplines the use of the instrument, and the dynamic path is shaped by the
interaction of voter patience, institutional gatekeeping, and the strength
of the moderate-defection channel.

% End of \section{Empirical Facts}
